% Worked solutions for Listas/lista-r4-moral-hazard-financing-capacity.md (15 exercises).
% Exercises 1, 3, 7, 8, 12, 13, 14, 15 are written here (label prefix r4e<n>-). Positions
% 2, 4, 5, 6, 9, 10, 11 are pasted from lista-topic4-moral-hazard-overhang-solutions.tex at the
% marker lines below (that file's labels keep their t4e<n>- prefix).
% Compile: pdflatex twice; check: pdffonts (Type 1 only, uni = yes).
% Every number is asserted in lista-r4-moral-hazard-financing-capacity-solutions_check.py.
\documentclass[11pt,a4paper]{article}
\usepackage[utf8]{inputenc}
\usepackage[T1]{fontenc}
\usepackage{lmodern}   % Type 1 vector fonts (replaces PK bitmaps)
\usepackage{cmap}      % ToUnicode CMap -> copyable text and accents
\usepackage[brazilian,english]{babel}
\usepackage{amsmath,amssymb,amsfonts}
\usepackage[a4paper,margin=2.2cm]{geometry}
\usepackage{booktabs}
\usepackage{enumitem}
\usepackage{fancyhdr}
\usepackage{xcolor}
\usepackage{microtype}
\usepackage{tikz}

\pagestyle{fancy}\fancyhf{}
\lhead{\footnotesize Finan\c{c}as II --- Corporativas (MPE Insper)}
\rhead{\footnotesize List R4 --- Worked solutions}
\cfoot{\footnotesize\thepage}
\renewcommand{\headrulewidth}{0.3pt}
\setlength{\headheight}{14pt}
\setlist{topsep=2pt,itemsep=1pt}

\newcommand{\exsec}[1]{\section*{\color{blue!55!black}#1}\addcontentsline{toc}{section}{#1}}
\newcommand{\notation}{\par\noindent\textbf{Notation.}}
\newcommand{\step}[2]{\par\medskip\noindent\textbf{Step #1 --- #2.}~}
\newcommand{\intuition}{\par\medskip\noindent\textbf{Intuition.}~}
\newcommand{\bookans}{\par\medskip\noindent\textbf{Book answer.}~}
\newcommand{\refbib}{\par\smallskip\noindent\textbf{Ref:}~}
\newcommand{\assume}{\par\smallskip\noindent\textbf{Assumption.}~}
\newcommand{\pt}{\foreignlanguage{brazilian}}
\newcommand{\asked}{\par\medskip\noindent\textbf{What is asked.}~}
\newcommand{\approach}{\par\smallskip\noindent\textbf{Approach.}~}
\newcommand{\answer}[1]{\par\smallskip\noindent\fbox{\textbf{Answer #1.}}~}
\newcommand{\pitfalls}{\par\medskip\noindent\textbf{Common mistakes.}}

\usepackage[hyphens]{url}
\begin{document}
\begin{center}{\Large\bfseries List R4 --- Moral hazard, financing capacity, debt vs.\ equity}\\[4pt]{\large Worked solutions}\end{center}

\noindent Statements are in \url{Listas/lista-r4-moral-hazard-financing-capacity.md}; each exercise is cited by number and printed page and is not restated. Exercises 2, 4, 5, 6, 9, 10 and 11 are the same book exercises as Topic~4 list exercises 4, 3, 5, 8, 9, 10 and 11; their solutions are taken from \url{lista-topic4-moral-hazard-overhang-solutions.tex}. Editions are the local copies: Tirole, \emph{The Theory of Corporate Finance} (2005); Berk = Berk, DeMarzo \& Harford, \emph{Fundamentals of Corporate Finance}, 2nd ed. (2011); Grinblatt--Titman, \emph{Financial Markets and Corporate Strategy}, 2nd ed. (2002); Ross = Ross, Westerfield, Jordan \& Lamb, \emph{Fundamentos de Administra\c{c}\~ao Financeira}, 9th ed. (2013). Unless stated otherwise, everybody is risk neutral and the riskless rate is $0$, as in Tirole.

\tableofcontents

% =============================================================================
\exsec{1. Agency relationship in the corporation --- Ross Cap.~1, Quest\~ao 7, p.~22}

\asked Four linked questions. (a) Who owns a corporation. (b) Through which process the owners
control the people who run it. (c) The main reason why the corporate form creates an agency
relationship. (d) Which problems can arise from that relationship.

\approach The answer comes from Ross's description of the corporation (\S1.2) and of the agency
problem (\S1.4): a principal (the shareholders) delegates decisions to an agent (management)
whose goals differ and whose actions the principal cannot fully observe. The same structure is
the entrepreneur--investor model of this list (Tirole \S3.2): the insider can take a private
benefit $B$ at the outsiders' expense, and the cure is to give her a stake, which limits what
can be promised to outsiders. Linking the two shows why the qualitative answer matters for
financing capacity.

\notation
\begin{itemize}
  \item Principal: the shareholders (owners); agent: the managers (executive officers) hired to run the firm.
  \item Agency problem: conflict of interest between principal and agent; agency costs: the costs of that conflict.
  \item For the formal link (Step 5): $p_H>p_L$, success probabilities when the manager behaves / misbehaves; $\Delta p=p_H-p_L$; $B$: private benefit from misbehaving; $R$: income on success; $R_b$: manager's reward on success.
\end{itemize}

\step{1}{(a) Who owns the corporation} We start from the legal structure, because ownership
determines who is the principal. The owners are the shareholders: they hold the residual claim
(what is left after all other claimants are paid) and the votes. Ross: \pt{``Em uma grande
empresa, acionistas e administradores formam grupos separados. Os acionistas elegem o conselho
de administra\c{c}\~ao, que ent\~ao seleciona os diretores executivos, encarregados da
administra\c{c}\~ao da empresa''} (p.~7).

\answer{(a)} The shareholders own the corporation.

\step{2}{(b) How the owners control management} We list the channels in the order Ross gives
them, from the formal chain of command to the market mechanisms that back it up.
\begin{enumerate}
  \item \emph{Board of directors.} Shareholders elect the board; the board hires, fires and
        sets the pay of the executives (\pt{``os acionistas controlam a empresa porque eles
        elegem os conselheiros''}, p.~7).
  \item \emph{Compensation.} Pay linked to performance and to the share price (stock options),
        so managers gain when the share price rises (pp.~13--14).
  \item \emph{Managerial labour market.} Managers who deliver value are promoted and paid more
        (p.~14).
  \item \emph{Proxy fights and takeovers.} Dissatisfied shareholders can replace the board
        (\pt{``disputa por procura\c{c}\~oes de voto''}), or another firm can buy the company
        and replace management (pp.~14--15). Ross notes that in Brazil, where control blocks are
        common, proxy fights are rarely effective.
\end{enumerate}

\answer{(b)} Shareholders elect the board, which appoints and pays the managers; this is
reinforced by performance-linked pay, the managerial labour market, and the threat of proxy
fights and takeovers.

\step{3}{(c) Why the corporate form creates an agency relationship} The reason is the
separation of ownership and control. Ownership is spread over many shareholders, none of whom
runs the firm, so management controls day-to-day decisions: \pt{``Essa dispers\~ao de
propriedade, sem d\'uvida, significa que quem realmente controla a empresa \'e a sua
administra\c{c}\~ao''} (p.~12). Ross defines the result: \pt{``O relacionamento entre
acionistas e administradores \'e chamado de relacionamento de ag\^encia''} (p.~12).

\answer{(c)} The separation between ownership (dispersed shareholders) and control
(professional managers) is the main reason.

\step{4}{(d) Which problems can arise} The agent may pursue her own goals when the principal
cannot observe or contract on what she does. Ross's examples, grouped:
\begin{enumerate}
  \item \emph{Different objectives}: rejecting a risky but value-increasing project to protect
        one's job; empire building (overpaying for acquisitions to run a bigger firm);
        emphasis on organisational survival and independence (pp.~12--13).
  \item \emph{Direct agency costs}: spending that benefits managers at shareholders' expense
        (Ross's example: an unnecessary corporate jet) and the cost of monitoring managers
        (external auditors) (p.~13).
  \item \emph{Indirect agency costs}: lost opportunities, such as the forgone project in item 1
        (p.~13).
  \item \emph{Brazilian context}: with controlling shareholders, the main conflict is between
        controllers and minority shareholders (private benefits of control, related-party
        transactions) (p.~13).
\end{enumerate}

\answer{(d)} Conflicts of interest (agency problems): perquisites, empire building,
excessive caution, entrenchment; they generate direct costs (perks, monitoring) and indirect
costs (forgone value). With concentrated ownership the conflict is controller versus minority.

\step{5}{Link to financing capacity} We now show, in the model used in the rest of this list,
why the agency problem limits what outsiders can be promised. If the manager behaves, success
has probability $p_H$; if she misbehaves, $p_L$ and she gets $B$. She behaves only if
\begin{align}
  p_H R_b &\ge p_L R_b+B \nonumber\\
  (p_H-p_L)R_b &\ge B \nonumber\\
  R_b &\ge \frac{B}{\Delta p}. \label{r4e1-IC}
\end{align}
With Tirole's illustration used in Exercise~5 ($R=150$, $B=30$, $\Delta p=0.4$),
\begin{align}
  \frac{B}{\Delta p} &= \frac{(30)}{(0.4)} = 75, \nonumber\\
  \frac{75}{R} &= \frac{75}{(150)} = 0.50, \nonumber
\end{align}
so half of the success income must stay with the insider and cannot be pledged to outside
owners or lenders. This is the ``pledgeable income'' $p_H(R-B/\Delta p)$ that drives
Exercises~5--15.

\intuition The agency problem is not about dishonest managers; it follows from three features
together: owners delegate, the agent's goals differ, and the agent's actions are hard to
observe. The control devices (board, pay, takeovers) align incentives only partly and at a
cost. In the financing models of this list, the same logic says the insider must keep a stake
of at least $B/\Delta p$ in success, so outsiders can be promised only what is left; that is
why net worth and monitoring matter for how much a firm can raise.
\pitfalls
\begin{itemize}
  \item Saying the managers or the board own the firm. The board represents the owners; the
        owners are the shareholders.
  \item Explaining the agency problem only by managerial bad faith. The cause is the
        separation of ownership and control combined with unobservable actions and different goals.
  \item Forgetting that monitoring costs (auditing, boards) are themselves agency costs, and
        that in Brazil the dominant conflict is controller versus minority shareholder.
\end{itemize}
\bookans The book gives no answer to this question.
\refbib{Ross Cap.~1, \S1.2 ``Formas de organiza\c{c}\~ao de neg\'ocios'' (sociedade por a\c{c}\~oes), pp.~7--8; \S1.4 ``O problema de ag\^encia e o controle da empresa'', pp.~12--17. Tirole Ch.~3, \S3.2.2, pp.~116--117.}

%% BEGIN-FROM lista-topic4-moral-hazard-overhang-solutions.tex EX 4 AS 2
\exsec{2. Short-term debt and incentive problems --- Grinblatt \& Titman 16.11, p.~593}

\asked A single conceptual question: why the conflicts between shareholders and creditors
(asset substitution, debt overhang, under-maintenance) do less damage when the firm finances itself
with debt that matures soon and must be rolled over, rather than with long-term debt.

\approach The conflicts exist because equity chooses the firm's strategy \emph{after} the terms of
the debt are fixed, so it can change the value of the creditors' claim (the lecture's asset
substitution and debt overhang). We write equity as firm value minus debt value and ask what
happens to the debt term when the debt is repriced after the choice is made: if new lenders price
the debt knowing the strategy, the debt term no longer depends on the strategy and equity
internalizes the full value consequences.

\notation
\begin{itemize}
  \item $j$: investment strategy (project choice, risk level, maintenance), chosen by the equity holders after the debt is in place; $X_j\ge 0$: resulting terminal firm value, $V_j=\mathbb{E}[X_j]$ (risk neutrality, $r_f=0$).
  \item $F$: face value of the debt; $L$: amount the firm owes when its short-term debt is rolled over; $D_j$, $E_j$: values of debt and equity under $j$.
\end{itemize}

\step{1}{Where the conflict comes from} We locate the conflict in the debt term of equity's value.
With long-term debt, $F$ is fixed before $j$ is chosen, and since $\max\{X-F,0\}=X-\min\{X,F\}$ in
every state,
\begin{align}
  E_j &= \mathbb{E}[\max\{X_j-F,0\}] \label{t4e4-E}\\
      &= V_j-D_j, \qquad D_j=\mathbb{E}[\min\{X_j,F\}]. \nonumber
\end{align}
Equity can gain by picking a $j$ that lowers $D_j$ (more risk, less investment, payouts) even
when $V_j$ falls: comparing two strategies, $E_j-E_k=(V_j-V_k)-(D_j-D_k)$, and the second term
is a transfer from creditors whenever $D_j<D_k$.

\step{2}{Debt that is repriced when the choice is made} We now let the debt be repriced after the
choice and show the transfer term disappears. If the debt matures before or at the
decision and must be rolled over, the new lenders see $j$ and set the new face $F_j$ so that
their claim is worth what they lend, $\mathbb{E}[\min\{X_j,F_j\}]=L$. Substituting into
\eqref{t4e4-E},
\begin{align}
  E_j &= V_j-(L), \label{t4e4-Est}
\end{align}
and $L$ does not depend on $j$: equity bears the full consequences of its choice and picks
the value-maximizing strategy. Short-term debt is worth close to its face whatever the firm
does, so there is little value to expropriate.

\step{3}{Illustration (Grinblatt \& Titman 16.1)} Take the firm of G\&T Exercise~16.1: \$100 of
debt, projects A (pays $120$ or $60$) and B (pays $101$ in both states), each state with
probability $0.5$, $r_f=0$. With long-term debt, equity is worth $0.5(20)+0.5(0)=10$ under A and
$0.5(1)+0.5(1)=1$ under B, so it picks the value-destroying A. Suppose instead the \$100 had to be
refinanced after the project is chosen. Firm values are
$V_A=0.5(120)+0.5(60)=90$ and $V_B=0.5(101)+0.5(101)=101$; by \eqref{t4e4-Est} with $L=100$,
\begin{align}
  E_A &= (90)-(100) = -10, & E_B &= (101)-(100) = 1. \nonumber
\end{align}
Under A the firm cannot raise \$100 (no claim on payoffs worth $90$ is worth $100$), so the
equity holders choose B, the positive-NPV project, instead of A.

\step{4}{Debt overhang} Myers (1977): if the existing debt matures before the firm must raise
money for a new project, the debt is repriced together with the project decision. The
project's benefit to creditors (a safer firm) then shows up as a lower borrowing rate that
accrues to equity, and positive-NPV projects are taken. In Unitron's case, the riskless
12\% project would lower the firm's 13\% borrowing rate and become attractive.

\step{5}{Asset substitution} The lending rate is renegotiated after the risk choice is made and
seen, so a risky choice raises the firm's own cost of funds instead of transferring wealth
from existing lenders. The threat of not being rolled over also disciplines the choice (lecture
Topic~4B, slide~16).

\step{6}{The cost} Rolling over exposes the firm to refinancing and interest-rate risk: an
unexpected rise in rates can push a highly levered firm into distress (hence the advice to
pair short-term debt with interest-rate hedges).

\answer{(16.11)} Short-term debt is repriced (rolled over) after the shareholders' investment and
risk choices are made and observed, so its terms reflect those choices: equity bears the full value
consequences, $E_j=V_j-L$, and cannot shift wealth from creditors through risk-taking or
under-investment. Short-term debt is also close to riskless, so there is little value to
expropriate. The price is rollover and interest-rate risk.

\intuition The conflicts exist because equity can change the value of a claim whose terms are
already fixed. The shorter the maturity, the sooner those terms are reset at prices that
reflect what equity has done, and the less wealth equity can take from creditors. The
incentive problem shrinks, and the firm takes on rollover risk in exchange.
\pitfalls
\begin{itemize}
  \item Saying short-term debt helps ``because it is safer''. The point is timing: the debt is priced
        after the decision, so the decision's effect on the creditors comes back to equity through the
        new rate.
  \item Confusing maturity with seniority. Seniority decides who is paid first in default; maturity
        decides when the terms are reset.
  \item Omitting the cost. Short maturities expose the firm to refinancing and interest-rate risk, so
        they are not a free cure.
\end{itemize}
\refbib{Grinblatt \& Titman Ch.~16, \S~16.4 ``The Use of Short-Term versus Long-Term Debt'', pp.~584--585 (Result~16.12). Lecture Topic~4B, slide~16.}
%% END-FROM

% =============================================================================
\exsec{3. Risk-increasing investment and leverage --- Berk 16.24, p.~494}

\asked Dynron (assets worth \$150 million) can sell \$50 million of pipeline and invest it in
a fibre-optic network that raises expected profits but makes the firm much riskier. Is this
switch more or less attractive to the shareholders if Dynron has debt than if it has none?

\approach This is the lecture's \emph{asset substitution} (Berk: ``excessive risk-taking'')
problem. Write equity as firm value minus debt, $E=V-D$; the change in equity from any decision
is then the project's NPV minus the change in the value of the debt. Levered equity is a call
option on the firm's assets and risky debt is riskless debt minus a put, so a riskier asset mix
raises the equity's value and lowers the debt's, holding the expected payoff fixed. Because the
statement gives no payoff distribution, the general argument is followed by an illustrative
two-state example (assumption flagged).

\notation
\begin{itemize}
  \item $X$: firm's total payoff next period (random); $V=\mathbb{E}[X]$: firm value (risk-neutral pricing, $r_f=0$ in the illustration).
  \item $F$: face value of the debt; $D=\mathbb{E}[\min\{X,F\}]$, $E=\mathbb{E}[\max\{X-F,0\}]$: values of debt and equity.
  \item $\mathit{NPV}$: change in $V$ caused by the switch; $\Delta E$, $\Delta D$: changes in $E$ and $D$.
\end{itemize}

\step{1}{Split the shareholders' gain} We write the shareholders' gain in a form that separates
value creation from transfers. Since $X=\min\{X,F\}+\max\{X-F,0\}$ in every state, $V=D+E$,
and therefore
\begin{align}
  \Delta E &= \Delta V-\Delta D \nonumber\\
           &= \mathit{NPV}-\Delta D. \label{r4e3-split}
\end{align}

\step{2}{No debt} With $F=0$, $D=0$ before and after, so $\Delta D=0$ and, by \eqref{r4e3-split},
$\Delta E=\mathit{NPV}$. An all-equity firm takes the switch only if it creates value.

\step{3}{Debt and risk} With risky debt, the payoff $\min\{X,F\}$ is a concave function of $X$
(it is $X$ below $F$ and flat at $F$ above). By Jensen's inequality, spreading out $X$ while
keeping its mean lowers $\mathbb{E}[\min\{X,F\}]$: the debt loses value, $\Delta D<0$, and from
\eqref{r4e3-split}
\begin{align}
  \Delta E &= \mathit{NPV}-\Delta D > \mathit{NPV}. \nonumber
\end{align}
Equivalently, $D=F-\mathbb{E}[\max\{F-X,0\}]$ (riskless debt minus a put), and the put gains
value with risk. If the debt stays riskless in every state ($X\ge F$ always), $\Delta D=0$ and
leverage changes nothing.

\assume Illustration (not in the statement): $r_f=0$, risk neutrality; the remaining pipeline
pays $100$ for sure; the fibre network pays $110$ or $0$ with probability $0.5$ each; before the
switch the firm pays $150$ for sure; debt face $F=120$. Amounts in \$ million.

\step{4}{Illustration: value created and value transferred} We compute the NPV, then each
claim before and after. The fibre network's expected payoff and NPV are
\begin{align}
  \mathbb{E}[\text{fibre}] &= (0.5)(110)+(0.5)(0) = 55, \nonumber\\
  \mathit{NPV} &= (55)-(50) = 5. \nonumber
\end{align}
After the switch the firm pays $100+110=210$ or $100+0=100$, so $V=(0.5)(210)+(0.5)(100)=155$.
Before the switch, $D=\min\{150,120\}=120$ and $E=150-120=30$. After it,
\begin{align}
  D &= (0.5)\min\{210,120\}+(0.5)\min\{100,120\} \nonumber\\
    &= (0.5)(120)+(0.5)(100) = 110, \nonumber\\
  E &= (0.5)\max\{210-120,0\}+(0.5)\max\{100-120,0\} \nonumber\\
    &= (0.5)(90)+(0.5)(0) = 45. \nonumber
\end{align}
So $\Delta E=45-30=15$ and $\Delta D=110-120=-10$; check with \eqref{r4e3-split}:
$15=(5)-(-10)$. Without debt the shareholders would gain only $155-150=5$.

\step{5}{Illustration: a value-destroying switch} Let the fibre network pay $90$ or $0$ instead.
Then $\mathbb{E}[\text{fibre}]=(0.5)(90)=45$ and $\mathit{NPV}=45-50=-5$; an unlevered firm
rejects it ($V$ falls from $150$ to $(0.5)(190)+(0.5)(100)=145$). With $F=120$,
\begin{align}
  D &= (0.5)\min\{190,120\}+(0.5)\min\{100,120\} = (0.5)(120)+(0.5)(100) = 110, \nonumber\\
  E &= (0.5)\max\{190-120,0\}+(0.5)(0) = (0.5)(70) = 35 > 30, \nonumber
\end{align}
so the levered shareholders accept it: $\Delta E=35-30=5=(-5)-(-10)$.

\step{6}{Boundary: riskless debt} With $F=80$, the debt is paid in full in both states after
the switch ($100\ge 80$): $D=(0.5)(80)+(0.5)(80)=80$, $\Delta D=0$, and
$\Delta E=(0.5)(130)+(0.5)(20)-(150-80)=75-70=5=\mathit{NPV}$, exactly as without debt.

\answer{} More attractive. With risky debt, the shareholders gain the NPV plus the value
transferred from the debt holders ($\Delta E=\mathit{NPV}-\Delta D$ with $\Delta D<0$); in the
illustration $15$ instead of $5$, and even a $-5$ NPV switch raises equity by $5$.

\intuition Limited liability caps the shareholders' downside at zero, so when the firm is
levered they keep the upside of a risky bet while the creditors absorb part of the downside.
That is why a levered firm likes risk more than an unlevered one, possibly to the point of
taking negative-NPV gambles. Creditors anticipate it and charge for it up front, so the
shareholders bear the agency cost ex ante, which is one reason for covenants and for the
financing limits studied in the Tirole exercises of this list.
\pitfalls
\begin{itemize}
  \item Answering ``less attractive, because the riskier firm has a higher cost of equity''.
        The cost of capital adjusts to risk; the effect that matters is the transfer from
        debt to equity.
  \item Ignoring that the result needs \emph{risky} debt: if the debt is paid in full in every
        state, leverage leaves the decision unchanged (Step~6).
  \item Treating the higher equity value as value creation. Firm value rises only by the NPV;
        the rest comes from the creditors.
\end{itemize}
\bookans The book gives no answer to this problem.
\refbib{Berk Ch.~16, \S16.6 ``Agency Costs and Information'', ``Equity-Debt Holder Conflicts'' and ``excessive risk-taking'', pp.~483--484. Lecture Topic~4B, slides 2--8.}

%% BEGIN-FROM lista-topic4-moral-hazard-overhang-solutions.tex EX 3 AS 4
\exsec{4. Debt overhang on vacant land --- Berk 16.25, pp.~494--495}

\asked (a) Today's values of equity and debt if the land stays vacant. (b) The NPV of developing
the land, for the firm as a whole. (c) Today's values of equity and debt if the shareholders put in
\$20 million and the land is developed. (d) Whether the shareholders gain from putting in that
money, given (c).

\approach All cash flows are riskless, so every claim is its payoff in one year discounted at
$r_f=10\%$; debt gets $\min\{\text{land value},15\}$ and equity the rest. This is the lecture's
\emph{debt overhang} (under-investment) problem: the debt is under water, so part of the value a
new project creates goes to the creditors, who pay none of its cost. Comparing the NPV (b) with the
shareholders' gain (d) shows the wedge.

\notation
\begin{itemize}
  \item $F=15$: face value of the debt, due in one year; $r_f=0.10$; all cash flows riskless; no taxes.
  \item $L_0=10$: value of the vacant land in one year; $L_1=35$: value of the developed land in one year; $K=20$: development cost paid today.
  \item $E$, $D$: values today of equity and debt.
\end{itemize}

\step{1}{(a) Land left vacant} We split the land's value in one year between the claims and
discount. In one year the firm is worth $10<15$: the creditors take $10$ and the shareholders get
$0$. Discounting at $r_f$,
\begin{align}
  D &= \frac{\min\{L_0,F\}}{1+r_f} \label{t4e3-Da}\\
    &= \frac{\min\{10,15\}}{1+0.10} \nonumber\\
    &= \frac{(10)}{(1.10)} \nonumber\\
    &= 9.0909, \nonumber\\
  E &= \frac{\max\{L_0-F,0\}}{1+r_f} = \frac{\max\{10-15,0\}}{1.10} = \frac{(0)}{(1.10)} = 0. \nonumber
\end{align}

\answer{(a)} Equity is worth \$0; debt is worth \$9.09 million.

\step{2}{(b) NPV of developing} The NPV counts only the cash flows that change with the decision.
Developing replaces land worth $10$ by land worth $35$, so the incremental cash flow in one year is
$35-10=25$, against $20$ paid today:
\begin{align}
  \mathit{NPV} &= -K+\frac{L_1-L_0}{1+r_f} \label{t4e3-NPV}\\
      &= -(20)+\frac{(35)-(10)}{(1.10)} \nonumber\\
      &= -20+\frac{25}{1.10} \nonumber\\
      &= -20+22.7273 \nonumber\\
      &= 2.7273. \nonumber
\end{align}
(Ignoring the vacant-land value forgone would give $-20+35/1.10=11.8182$, the NPV measured
against a firm with no land at all; the decision-relevant figure is \eqref{t4e3-NPV}.)

\answer{(b)} $\mathit{NPV}=\$2.73$ million $>0$: developing creates value.

\step{3}{(c) Values if developed with \$20 of new equity} We repeat Step 1 with the developed land.
The \$20 is spent on development, so it is not an asset in one year. In one year the firm is worth
$35\ge 15$: the debt is paid in full, $\min\{35,15\}=15$, and equity keeps $\max\{35-15,0\}=20$. Hence
\begin{align}
  D &= \frac{(15)}{(1.10)} \label{t4e3-Dc}\\
    &= 13.6364, \nonumber\\
  E &= \frac{(35)-(15)}{(1.10)} \label{t4e3-Ec}\\
    &= \frac{20}{1.10} \nonumber\\
    &= 18.1818. \nonumber
\end{align}

\answer{(c)} Equity is worth \$18.18 million and debt \$13.64 million today.

\step{4}{(d) Will the shareholders put up \$20?} The shareholders compare what they hold after
contributing with what they hold if they do nothing. Their gain is the equity value
\eqref{t4e3-Ec}, minus the equity value without development (Step 1), minus the contribution:
\begin{align}
  \Delta E &= E_{\text{dev}}-E_{\text{vacant}}-K \nonumber\\
           &= (18.1818)-(0)-(20) \nonumber\\
           &= -1.8182 < 0, \nonumber\\
  \Delta D &= D_{\text{dev}}-D_{\text{vacant}} \nonumber\\
           &= (13.6364)-(9.0909) \nonumber\\
           &= 4.5455. \nonumber
\end{align}
No: they would lose \$1.82 million. The sum $-1.8182+4.5455=2.7273$ is the NPV
\eqref{t4e3-NPV}; the creditors capture more than all of it.

\answer{(d)} No. The shareholders would pay \$20 million for an equity stake worth \$18.18 million
(a loss of \$1.82 million), while the debt gains \$4.55 million; the positive-NPV project is passed up.

\intuition The debt is under water ($10<15$), so the first \$5 million of any value the
project creates goes to the creditors, who pay none of its cost. The shareholders bear the
\$20 million and receive only the value above the \$15 million face. This is Myers' debt
overhang: leverage acts as a tax on new investment, and a positive-NPV project is passed up.
Writing the debt down (renegotiation, as in the lecture) or senior new financing would fix it.
\pitfalls
\begin{itemize}
  \item Computing the NPV as $-20+35/1.10=11.82$: the vacant land worth \$10 million in one year is
        given up, so it is part of the project's cost.
  \item Answering (d) with $E_{\text{dev}}=18.18>0$, ``so yes''. The shareholders must compare the
        stake with the \$20 million they put in.
  \item Valuing the debt in (a) at the face value $15/1.10$; the creditors can get at most what the
        land is worth, $\min\{10,15\}$.
\end{itemize}
\refbib{Berk Ch.~16, \S~16.6 ``Agency Costs and Information'', ``Equity-Debt Holder Conflicts'' (under-investment), pp.~483--484. Lecture Topic~4B, slides 19--30.}
%% END-FROM

%% BEGIN-FROM lista-topic4-moral-hazard-overhang-solutions.tex EX 5 AS 5
\exsec{5. Random financing --- Tirole 3.1, p.~144}

\asked (i) Write the condition under which investors are willing to sign a contract in which the
borrower contributes $\hat A$ of her cash, the project then goes ahead only with probability $x$,
and she gets $R_b$ on success. (ii) Show that a borrower with too little cash for certain financing
($A<\bar A$) should put \emph{all} her cash in, and describe how the financing probability depends
on $A$.

\approach This is the Holmstr\"om--Tirole fixed-investment model of the lecture (Topic~4A). The two
constraints are the borrower's incentive constraint (IC$_b$), which caps the \emph{pledgeable income}
at $\mathcal P=p_H(R-B/\Delta p)$, and the investors' breakeven constraint (IR$_l$). Below $\bar A$
the project cannot be financed for sure; a lottery over financing scales the investors' outlay and
the pledgeable income by $x$ but not the borrower's contribution, which convexifies the
indivisible project. We write (IR$_l$) for the lottery, show (IC$_b$) must hold, and maximise the
borrower's utility over $(x,\hat A)$.

\notation
\begin{itemize}
  \item $I>0$: fixed investment; $A\in(0,\bar A)$: borrower's cash; $\hat A\in[0,A]$: the part she contributes; $x\in[0,1]$: probability that the project is financed.
  \item $R>0$: income on success (0 on failure); $R_b\ge 0$: borrower's reward on success; $R_l=R-R_b$: investors' claim on success.
  \item $p_H>p_L$: success probabilities when she behaves / misbehaves; $\Delta p=p_H-p_L$; $B>0$: private benefit from misbehaving.
  \item $\mathcal P=p_H(R-B/\Delta p)$: pledgeable income; $\bar A=I-\mathcal P$: minimum cash for financing with probability~1 (the statement's definition); $U_b$: borrower's net utility. Rate of interest 0, everybody risk neutral.
\end{itemize}

\assume The borrower hands over $\hat A$ when the contract is signed, whether or not the
lottery then selects the project (this is the reading in Tirole's answer). If $\hat A$ were
paid only when the project goes ahead, the expected outlay would be $x(I-\hat A)$ and the
breakeven condition $x(I-\hat A)\le x\mathcal P$ would require $\hat A\ge\bar A>A$: randomizing
would achieve nothing.

\step{1}{(i) Investors' breakeven condition} Investors receive $\hat A$ for sure, pay $I$ with
probability $x$, and, if the borrower behaves, collect $R_l$ with probability $p_H$ when the
project runs:
\begin{align}
  \underbrace{\hat A}_{\text{received}}+\underbrace{x\,p_H(R-R_b)}_{\text{expected repayment}}
     &\ge \underbrace{x I}_{\text{expected outlay}} \nonumber\\
  x\,p_H(R-R_b) &\ge x I-\hat A. \label{t4e5-IR}
\end{align}

\answer{(i)} $x\,p_H(R-R_b)\ge xI-\hat A$ (with $R_b\ge B/\Delta p$ so that she behaves).

\step{2}{The contract must induce good behaviour} We check that a contract under which she shirks
is never worth signing, so (IC$_b$) can be imposed. Suppose $(\Delta p)R_b<B$, so she shirks.
Investors break even only if $x p_L(R-R_b)\ge xI-\hat A$, i.e.\ $-\hat A\le x p_L(R-R_b)-xI$.
Her net utility is then
\begin{align}
  x(p_L R_b+B)-\hat A &\le x(p_L R_b+B)+\big(x p_L(R-R_b)-xI\big) \nonumber\\
                      &= x p_L R_b+xB+x p_L R-x p_L R_b-xI \nonumber\\
                      &= x(p_L R+B-I) < 0 \label{t4e5-shirk}
\end{align}
by (3.2), worse than not signing. So the contract satisfies (IC$_b$), $(\Delta p)R_b\ge B$, i.e.\
$R_b\ge B/\Delta p$.

\step{3}{Breakeven in terms of $\bar A$} We combine the two constraints into one feasibility
condition on $(x,\hat A)$. (IC$_b$) gives $R-R_b\le R-B/\Delta p$, so the
left side of \eqref{t4e5-IR} is at most $x\,p_H(R-B/\Delta p)=x\mathcal P$. A contract exists for $(x,\hat A)$ only if
\begin{align}
  x I-\hat A &\le x\,p_H\left(R-\frac{B}{\Delta p}\right) \nonumber\\
  x I-\hat A &\le x(\mathcal P) \nonumber\\
  x(I-\mathcal P) &\le \hat A \nonumber\\
  x(\bar A) &\le \hat A \nonumber\\
  x &\le \frac{\hat A}{\bar A}, \label{t4e5-x}
\end{align}
substituting $I-\mathcal P=\bar A$ from the statement's definition (third to fourth line) and
dividing by $\bar A>0$ (last line). Conversely, if \eqref{t4e5-x} holds, $R_b=B/\Delta p$ satisfies
both constraints, so \eqref{t4e5-x} is necessary and sufficient.

\step{4}{(ii) The borrower's objective} We write her utility and use the binding breakeven
constraint to express it in terms of $x$ only. She keeps $A-\hat A$ and, with probability $x$, earns
$R_b$ with probability $p_H$; net of the $A$ she would consume anyway,
\begin{align}
  U_b &= x\,p_H R_b+(A-\hat A)-A \nonumber\\
      &= x\,p_H R_b-\hat A. \label{t4e5-Ub}
\end{align}
Competitive investors make \eqref{t4e5-IR} bind; expanding and solving for $x p_H R_b$,
\begin{align}
  x p_H R-x p_H R_b &= xI-\hat A \nonumber\\
  x p_H R_b &= x p_H R-xI+\hat A. \nonumber
\end{align}
Substituting into \eqref{t4e5-Ub},
\begin{align}
  U_b &= (x p_H R-xI+\hat A)-\hat A \nonumber\\
      &= x\,(p_H R-I). \label{t4e5-Ubx}
\end{align}

\step{5}{(ii) Optimal $\hat A$ and $x$} We maximise \eqref{t4e5-Ubx} over the feasible set
\eqref{t4e5-x}. Differentiating,
\begin{align}
  \frac{\partial U_b}{\partial x} = p_H R-I>0, \qquad \frac{\partial U_b}{\partial \hat A}=0, \nonumber
\end{align}
so $U_b$ rises with $x$ and $\hat A$ matters only through feasibility. By \eqref{t4e5-x}, the
largest feasible $x$ is $\hat A/\bar A$ (below 1 because $\hat A\le A<\bar A$), which increases with
$\hat A$: substituting it, $U_b=(\hat A/\bar A)(p_HR-I)$, increasing in $\hat A$ on $[0,A]$. So
she contributes everything:
\begin{align}
  \hat A^* &= A, \nonumber\\
  x^* &= \frac{A}{\bar A}, \label{t4e5-xstar}\\
  U_b^* &= \frac{A}{\bar A}\,(p_H R-I). \label{t4e5-Ubstar}
\end{align}
At $x^*$ the binding \eqref{t4e5-IR}, divided by $x^*=A/\bar A$, gives
$p_H(R-R_b)=I-A/x^*=I-\bar A=\mathcal P$, so $R-R_b=R-B/\Delta p$ and $R_b=B/\Delta p$: (IC$_b$) binds.

\step{6}{(ii) How $x^*$ varies with $A$} $x^*=A/\bar A$ rises linearly from 0 at $A=0$ to 1 at
$A=\bar A$, and stays at 1 above $\bar A$. Utility \eqref{t4e5-Ubstar} therefore rises smoothly
to the full NPV instead of jumping at $\bar A$ as in (3.6).

\emph{Numerical illustration.} $I=100$, $R=150$, $p_H=0.8$, $p_L=0.4$, $B=30$, $A=10$: (3.2)
holds ($0.4(150)-100+30=-10<0$); $B/\Delta p=75$; $\mathcal P=0.8(150-75)=60$; $\bar A=100-60=40$;
NPV $=0.8(150)-100=20$. Then $x^*=10/40=0.25$ and $U_b^*=0.25(20)=5$; investors pay
$0.25(100)-10=15$ in expectation and receive $0.25(60)=15$.

\answer{(ii)} $\hat A^*=A$; the financing probability is $x^*=A/\bar A$, rising linearly from $0$
at $A=0$ to $1$ at $A=\bar A$ (and $x^*=1$ for $A\ge\bar A$); her utility is $x^*(p_HR-I)$.

\intuition The obstacle is not the size of the NPV but the pledgeable income per unit of
investment. A poor borrower can still buy a \emph{fraction} of the deal: each unit of her
cash covers the investors' shortfall on $1/\bar A$ units of probability of financing. Since
she captures the whole NPV whenever the project runs, she wants as high a probability as her
cash can buy and puts all of it in. Lotteries convexify the indivisible project.
\pitfalls
\begin{itemize}
  \item Writing the breakeven as $x p_H(R-R_b)\ge x(I-\hat A)$, i.e.\ taking her contribution only when
        the project goes ahead. Then $x$ cancels and randomisation achieves nothing.
  \item Using $R$ instead of $R-B/\Delta p$ as the investors' maximal claim: without (IC$_b$) she
        shirks, and by (3.2) no one breaks even.
  \item Claiming she keeps some cash ``for safety''. Cash kept aside earns nothing extra, while each unit
        put in raises $x$ by $1/\bar A$ and her utility by $(p_HR-I)/\bar A>0$.
\end{itemize}
\bookans Matches the book's answer (p.~571): breakeven $xI-A\le x p_H(R-R_b)$, equivalently
$x\bar A\le A$; $U_b=x(p_H R-I)$, so $x^*=A/\bar A$, rising from 0 to 1 as $A$ goes from 0
to $\bar A$.
\refbib{Tirole Ch.~3, \S~3.2.1--3.2.2 ``The Fixed-Investment Model'', ``The Lenders' Credit Analysis'', pp.~115--117 (footnote~9, p.~116). Lecture Topic~4A, slides 26--34.}
%% END-FROM

%% BEGIN-FROM lista-topic4-moral-hazard-overhang-solutions.tex EX 8 AS 6
\exsec{6. Project riskiness and credit rationing --- Tirole 3.15, p.~150}

\asked Two versions of the fixed-investment project have the same cost, the same expected income
under good behaviour and the same effect of shirking on the success probability, but B succeeds less
often with a larger payoff. Decide which version can be financed for a wider range of the
entrepreneur's cash $A$, i.e.\ is less prone to credit rationing.

\approach In the Holmstr\"om--Tirole fixed-investment model (lecture Topic~4A) a project is
financed iff $A\ge\bar A=I-\mathcal P$, where $\mathcal P=p_H(R-B/\Delta p)$ is the
\emph{pledgeable income}. So the variant with the lower threshold $\bar A$, i.e.\ the higher
pledgeable income, is less rationed. We write $\bar A$ as non-pledgeable rent minus NPV, use equal
NPVs, and compare the rents.

\notation
\begin{itemize}
  \item $j\in\{A,B\}$: project variant; $p_H^j>p_L^j$: success probabilities with good / bad behaviour; $R^j>0$: income on success (0 on failure).
  \item $\Delta p=p_H^A-p_L^A=p_H^B-p_L^B>0$: common; $B>0$: private benefit (common); $I$: common investment; $A$: borrower's cash.
  \item Maintained: $p_H^A R^A=p_H^B R^B$ (equal NPV) and $p_H^A>p_H^B$ (B riskier); (3.1), (3.2) hold for both variants.
  \item $\mathcal P^j=p_H^j(R^j-B/\Delta p)$: pledgeable income; $\bar A^j$: minimum cash for financing variant $j$ (3.3).
\end{itemize}

\step{1}{The minimum stake is the same} We first derive the borrower's minimum reward in each
variant. Variant $j$'s incentive constraint is $p_H^jR_b^j\ge p_L^jR_b^j+B$, i.e.\
$(p_H^j-p_L^j)R_b^j\ge B$, i.e.\ $R_b^j\ge B/\Delta p$: the same for both variants, because $B$ and
$\Delta p$ are common.

\step{2}{Financing thresholds} We compute each variant's threshold from (IR$_l$),
$\mathcal P^j\ge I-A$. From (3.3),
\begin{align}
  \bar A^j &= I-\mathcal P^j \nonumber\\
           &= I-p_H^j\left(R^j-\frac{B}{\Delta p}\right) \nonumber\\
           &= I-p_H^jR^j+p_H^j\frac{B}{\Delta p} \nonumber\\
           &= p_H^j\frac{B}{\Delta p}-(p_H^j R^j-I). \label{t4e8-Abar}
\end{align}
The first term is the non-pledgeable income: the expected rent the borrower must be left.

\step{3}{Compare} Subtracting \eqref{t4e8-Abar} for B from that for A, and substituting
$p_H^B R^B=p_H^A R^A$,
\begin{align}
  \bar A^A-\bar A^B &= \left[p_H^A\frac{B}{\Delta p}-(p_H^A R^A-I)\right]-\left[p_H^B\frac{B}{\Delta p}-((p_H^A R^A)-I)\right] \nonumber\\
                    &= p_H^A\frac{B}{\Delta p}-p_H^AR^A+I-p_H^B\frac{B}{\Delta p}+p_H^AR^A-I \nonumber\\
                    &= (p_H^A-p_H^B)\frac{B}{\Delta p} > 0, \label{t4e8-diff}
\end{align}
positive because $p_H^A>p_H^B$, $B>0$, $\Delta p>0$. The riskier variant B is financed for every
$A\ge\bar A^B$, including $A\in[\bar A^B,\bar A^A)$ where A is not: \textbf{variant B is less prone
to credit rationing}.

\answer{(3.15)} Variant B, the riskier one: $\bar A^B=\bar A^A-(p_H^A-p_H^B)B/\Delta p<\bar A^A$, so
every entrepreneur financed under A is financed under B, and those with
$A\in[\bar A^B,\bar A^A)$ only under B (given that both variants need good behaviour, (3.2)).

\step{4}{Numerical illustration} $I=100$, $\Delta p=0.2$, $B=9$; A: $p_H^A=0.8$, $p_L^A=0.6$,
$R^A=150$; B: $p_H^B=0.6$, $p_L^B=0.4$, $R^B=200$. Both NPVs are $120-100=20$; (3.2) holds
($0.6(150)-100+9=-1<0$ and $0.4(200)-100+9=-11<0$); $B/\Delta p=45$. Then
$\mathcal P^A=0.8(105)=84$, $\bar A^A=16$; $\mathcal P^B=0.6(155)=93$, $\bar A^B=7$; and
$\bar A^A-\bar A^B=(0.2)(45)=9$, as \eqref{t4e8-diff} says.

\step{5}{Beyond the maintained assumption (the book's extension)} If (3.2) fails for the safe
variant, a contract that leaves all income to investors ($R_b=0$, the borrower misbehaves)
may finance it: investors then receive $p_L^AR^A$, and the borrower, who gets the NPV with
shirking, $p_L^AR^A+B-I$, when investors break even, wants it if that NPV is positive. The safe
variant can then be financed and the risky one not when (in order) B cannot be financed with good
behaviour, B cannot be financed with shirking, A can be financed with shirking, and A's NPV with
shirking is positive:
\begin{align}
  I-A &> p_H^B\left(R^B-\frac{B}{\Delta p}\right), \quad I-A>p_L^B R^B, \quad I-A\le p_L^A R^A, \quad I<p_L^A R^A+B. \label{t4e8-ext}
\end{align}
These conditions are compatible. With $A=0$, $I=120$, $\Delta p=0.3$, $B=45$, $p_H^A=0.9$,
$R^A=200$, $p_H^B=0.5$, $R^B=360$ (both $p_H R=180$): $p_H^B(R^B-B/\Delta p)=0.5(210)=105<120$;
$p_L^B R^B=0.2(360)=72<120$; $p_L^A R^A=0.6(200)=120\ge 120$; $p_L^A R^A+B=165>120$. Here the
risky variant is rationed and the safe one is funded. Good behaviour adds more pledgeable
income when the success payoff is large, which favours the risky variant only when good
behaviour is needed for funding.

\intuition The borrower's rent $B/\Delta p$ is paid only on success. A safer project succeeds
more often, so it pays out the rent more often and leaves less expected income to pledge: the
non-pledgeable rent $p_H B/\Delta p$ is higher for the same NPV. Riskier projects with the same mean concentrate
income in success, where it is pledged anyway.
\pitfalls
\begin{itemize}
  \item Importing the Stiglitz--Weiss intuition that riskier borrowers are rationed first. Here
        riskiness lowers the rent $p_HB/\Delta p$ and so \emph{raises} pledgeable income.
  \item Computing the threshold as $I-p_HR$ (all income pledged), which is the same for both variants
        by $p_H^AR^A=p_H^BR^B$ and hides the difference.
  \item Forgetting that the comparison holds only while good behaviour is required for funding: if
        (3.2) fails for one variant, the book's extension (Step 5) can reverse the ranking.
\end{itemize}
\bookans Matches the book's answer (p.~578): with the same minimum reward $B/\Delta p$, the
safe project A is financed for a smaller range of $A$, i.e.\ it is more prone to credit
rationing; the book then shows, with conditions \eqref{t4e8-ext}, that the risky project may be
the more rationed one if high-powered incentives are not necessarily called for. One slip in
the book's prose: it says the third condition means that ``the risky project'' generates
enough pledgeable income when investors get all the cash flow, but that condition,
$I-A\le p_L^A R^A$, is about the safe variant A.
\refbib{Tirole Ch.~3, \S~3.2.1--3.2.2 ``The Fixed-Investment Model'', ``The Lenders' Credit Analysis'', pp.~115--117. Lecture Topic~4A, slides 26--35 (determinants of credit rationing).}
%% END-FROM

% =============================================================================
\exsec{7. Equity multiplier and active monitoring --- Tirole Exercise 3.8, p.~147}

\asked (i) In the variable-investment model, derive how much the borrower can invest per unit
of her own cash (the equity multiplier) and show it equals $1/(1-\rho_0)$. (ii) A monitor can be
hired who, at cost $cI$, cuts the borrower's private benefit to $b(c)I$. (ii-a) Write the
incentive constraints of the monitor and of the borrower. (ii-b) Derive the new multiplier.
(ii-c) Show which objective the borrower maximizes when choosing $c$.

\approach Holmstr\"om--Tirole with variable investment (Tirole \S3.4): the borrower's
incentive constraint fixes the minimum stake $BI/\Delta p$, so only $\rho_0=p_H(R-B/\Delta p)$
per unit of investment can be pledged; competitive investors break even, so the borrower gets
the whole NPV and invests as much as the pledgeable income allows. Monitoring (Tirole \S9.2.1)
lowers the stake the borrower needs, which raises pledgeable income, but the monitor must also be
paid a stake; with ``no shortage of monitoring capital'' the monitor pays for his stake up front,
so his cost is the only thing the borrower loses.

\notation
\begin{itemize}
  \item $I\ge 0$: investment; $A>0$: borrower's cash; income $RI$ on success, $0$ on failure.
  \item $p_H$, $p_L=p_H-\Delta p$: success probabilities if the borrower behaves / misbehaves; $BI$: private benefit from misbehaving.
  \item $\rho_1=p_HR>1$: expected income per unit; $\rho_0=p_H(R-B/\Delta p)<1$: pledgeable income per unit; $k$: equity multiplier, $I=kA$.
  \item $R_b$, $R_m$, $R_l$: rewards on success of the borrower, the monitor and the uninformed investors ($R_b+R_m+R_l=RI$); $0$ on failure (limited liability).
  \item $c\ge 0$: monitoring intensity, private cost $cI$ to the monitor; $b(c)$: private benefit per unit under monitoring, $b(0)=B$, $b'<0$; $I_m$: monitor's cash contribution.
  \item $U_b$: borrower's net utility (expected payoff minus the cash $A$ she could consume).
\end{itemize}

\step{1}{(i) The borrower's program} The borrower offers the contract; she maximizes her net
utility subject to her own incentive constraint and the investors' breakeven constraint:
\begin{align}
  \max_{I,\,R_b}\ & p_HR_b-A \label{r4e7-obj}\\
  \text{s.t. }\ & (\Delta p)R_b\ge BI, \tag{IC$_b$}\\
                & p_H(RI-R_b)\ge I-A. \tag{IR$_l$}
\end{align}

\step{2}{(i) The investors break even and the borrower gets the NPV} If (IR$_l$) were slack,
raising $R_b$ a little would keep it satisfied, relax (IC$_b$), and raise \eqref{r4e7-obj}; so
(IR$_l$) binds:
\begin{align}
  p_HRI-p_HR_b &= I-A \nonumber\\
  p_HR_b &= p_HRI-I+A. \label{r4e7-bind}
\end{align}
Substituting \eqref{r4e7-bind} into \eqref{r4e7-obj},
\begin{align}
  U_b &= (p_HRI-I+A)-A \nonumber\\
      &= (\rho_1-1)\,I, \label{r4e7-Ub}
\end{align}
which increases with $I$ because $\rho_1>1$: she invests as much as the constraints allow.

\step{3}{(i) The largest investment} Multiplying (IC$_b$) by $p_H/\Delta p>0$ gives
$p_HR_b\ge p_HBI/\Delta p$. Substituting $p_HR_b$ from \eqref{r4e7-bind},
\begin{align}
  \rho_1I-I+A &\ge \frac{p_HB}{\Delta p}I \nonumber\\
  A &\ge I\Big(1-\rho_1+\frac{p_HB}{\Delta p}\Big) \nonumber\\
  A &\ge I(1-\rho_0) \qquad\Big(\text{since }\rho_0=\rho_1-\frac{p_HB}{\Delta p}\Big) \nonumber\\
  I &\le \frac{A}{1-\rho_0} \qquad(\text{dividing by }1-\rho_0>0). \label{r4e7-k}
\end{align}
The bound is reached (both constraints bind), so $k=1/(1-\rho_0)$ and, by \eqref{r4e7-Ub},
$U_b=(\rho_1-1)A/(1-\rho_0)$.

\answer{(i)} $I=kA$ with $k=\dfrac{1}{1-\rho_0}=\dfrac{1}{1-p_H(R-B/\Delta p)}>1$.

\emph{Numerical illustration.} $p_H=0.8$, $p_L=0.4$ ($\Delta p=0.4$), $R=1.5$, $B=0.3$,
$A=10$. Then $\rho_1=(0.8)(1.5)=1.2$; $B/\Delta p=0.3/0.4=0.75$; $\rho_0=(0.8)(1.5-0.75)=0.6$;
the assumption $1>p_LR+B=(0.4)(1.5)+0.3=0.9$ holds. Hence $k=1/(1-0.6)=2.5$, $I=(2.5)(10)=25$,
$U_b=(1.2-1)(25)=5$.

\step{4}{(ii-a) The two incentive constraints} To induce monitoring $c$, the borrower's stake is
set so that she behaves when monitored, and the monitor's stake so that he monitors. Under
monitoring the borrower's private benefit is $b(c)I$, so
\begin{align}
  (\Delta p)R_b \ge b(c)I. \label{r4e7-ICb}
\end{align}
If the monitor skips monitoring he saves $cI$, but the borrower's private benefit is back to
$BI>b(c)I$; with $R_b$ set at the level \eqref{r4e7-ICb} she then misbehaves and success falls
to $p_L$. The monitor monitors if $p_HR_m-cI\ge p_LR_m$, that is,
\begin{align}
  (\Delta p)R_m \ge cI. \label{r4e7-ICm}
\end{align}

\answer{(ii-a)} $(\Delta p)R_m\ge cI$ (monitor) and $(\Delta p)R_b\ge b(c)I$ (borrower).

\step{5}{(ii-b) The monitor's contribution} With plentiful monitoring capital, the monitor
accepts any deal in which he breaks even, so the borrower makes him pay for his stake: his cash
contribution equals his expected income net of his cost. At the cheapest stake,
$R_m=cI/\Delta p$ from \eqref{r4e7-ICm}, and
\begin{align}
  I_m &= p_HR_m-cI \nonumber\\
      &= p_H\frac{cI}{\Delta p}-cI \nonumber\\
      &= \Big(\frac{p_H}{\Delta p}-1\Big)cI = \frac{p_L}{\Delta p}\,cI. \label{r4e7-Im}
\end{align}

\step{6}{(ii-b) The multiplier} The uninformed investors put up $I-A-I_m$ and receive $R_l$ on
success; by the argument of Step~2 their constraint binds, and by the argument of Step~3
(utility increasing in $I$) both incentive constraints bind:
$R_b=b(c)I/\Delta p$, $R_m=cI/\Delta p$. Their breakeven condition is
\begin{align}
  p_H(RI-R_b-R_m) &= I-A-I_m. \label{r4e7-IRl}
\end{align}
Substituting $R_b$, $R_m$ and \eqref{r4e7-Im},
\begin{align}
  p_H\Big(R-\frac{b(c)}{\Delta p}-\frac{c}{\Delta p}\Big)I &= I-A-p_H\frac{c}{\Delta p}I+cI \nonumber\\
  p_H\Big(R-\frac{b(c)}{\Delta p}\Big)I-p_H\frac{c}{\Delta p}I &= I-A-p_H\frac{c}{\Delta p}I+cI \nonumber\\
  p_H\Big(R-\frac{b(c)}{\Delta p}\Big)I &= I-A+cI \qquad(\text{the }p_Hc I/\Delta p\text{ terms cancel}) \nonumber\\
  A &= I\Big[1+c-p_H\Big(R-\frac{b(c)}{\Delta p}\Big)\Big]. \nonumber
\end{align}
Hence $I=k(c)A$ with
\begin{align}
  k(c) &= \frac{1}{1+c-p_H\big(R-b(c)/\Delta p\big)}. \label{r4e7-kc}
\end{align}
To write it with $\rho_0$, use $p_H(R-b/\Delta p)=p_H(R-B/\Delta p)+(p_H/\Delta p)(B-b)=\rho_0-(p_H/\Delta p)(b-B)$:
\begin{align}
  k(c) &= \frac{1}{1+c-\rho_0+\frac{p_H}{\Delta p}\big(b(c)-B\big)}
        = \frac{1}{1-\rho_0+c+\frac{p_H}{\Delta p}\big[b(c)-B\big]}. \label{r4e7-kc2}
\end{align}
At $c=0$, $b(0)=B$ and \eqref{r4e7-kc2} reduces to $1/(1-\rho_0)$. Monitoring raises the
multiplier if and only if $c+(p_H/\Delta p)[b(c)-B]<0$: the pledgeable income it frees,
$(p_H/\Delta p)[B-b(c)]$ per unit, exceeds its cost $c$.

\answer{(ii-b)} $k(c)=\dfrac{1}{1-\rho_0+c+(p_H/\Delta p)[b(c)-B]}$ (positive denominator assumed).

\step{7}{(ii-c) The borrower's objective} Neither the monitor nor the uninformed investors earn
a surplus, so the borrower should receive the NPV net of the monitoring cost. To confirm it,
solve \eqref{r4e7-IRl} for $p_HR_b$, using $I_m=p_HR_m-cI$:
\begin{align}
  p_HR_b &= p_HRI-p_HR_m-I+A+I_m \nonumber\\
         &= p_HRI-p_HR_m-I+A+p_HR_m-cI \nonumber\\
         &= (\rho_1-1-c)I+A, \nonumber\\
  U_b &= p_HR_b-A = (\rho_1-1-c)\,I = (\rho_1-1-c)\,k(c)\,A. \label{r4e7-Ubc}
\end{align}
Since $A$ is given, she chooses $c$ to
\begin{align}
  \max_c\ \frac{\rho_1-1-c}{1-\rho_0+c+(p_H/\Delta p)[b(c)-B]}. \label{r4e7-max}
\end{align}

\answer{(ii-c)} The borrower maximizes $(\rho_1-1-c)\,k(c)$, i.e.\ \eqref{r4e7-max}: the NPV
per unit of investment (net of monitoring cost) times the multiplier.

\step{8}{About the expression printed in the statement} The statement prints the denominator
$1-\rho_0+(p_H/\Delta p)[b(c)+c-B]$, which exceeds ours by $(p_H/\Delta p-1)c=(p_L/\Delta p)c=I_m/I$.
That is the multiplier when the monitor brings \emph{no} cash ($I_m=0$): then
$p_H(R-b/\Delta p-c/\Delta p)I=I-A$, i.e.\ $A=I[1-\rho_0+(p_H/\Delta p)(b+c-B)]$. But then the
monitor keeps a rent $p_HR_m-cI=(p_L/\Delta p)cI$ and the borrower's utility is
$(\rho_1-1-(p_H/\Delta p)c)I$, so the numerator would change too. Under the statement's own
assumption (the monitor breaks even), the consistent objective is \eqref{r4e7-max}, which is
also the book's answer (p.~574).

\step{9}{Choosing $c$: first-order condition and an example} Let $N(c)=\rho_1-1-c$ and
$D(c)=1-\rho_0+c+(p_H/\Delta p)[b(c)-B]$, with $N'=-1$ and $D'=1+(p_H/\Delta p)b'(c)$. By the
quotient rule,
\begin{align}
  \frac{d}{dc}\frac{N}{D} &= \frac{N'D-ND'}{D^2} = \frac{-D-N\big[1+(p_H/\Delta p)b'(c)\big]}{D^2}. \nonumber
\end{align}
Setting the numerator to zero,
\begin{align}
  -D-N-N\frac{p_H}{\Delta p}b'(c) &= 0 \nonumber\\
  -\frac{p_H}{\Delta p}b'(c) &= \frac{D+N}{N} = 1+\frac{D}{N}: \label{r4e7-foc}
\end{align}
the marginal gain in pledgeable income per unit, $-(p_H/\Delta p)b'(c)$, equals the marginal
cost, which hits both the NPV and the multiplier.

\emph{Example} (parameters of the illustration in (i), and $b(c)=B/(1+20c)$). Let $u=1+20c$, so
$c=(u-1)/20$, $b=0.3/u$, $b'(c)=-20B/u^2=-6/u^2$ and $p_H/\Delta p=2$. Then
\begin{align}
  N &= 0.2-\frac{u-1}{20} = \frac{5-u}{20}, \nonumber\\
  D &= 0.4+\frac{u-1}{20}+2\Big(\frac{0.3}{u}-0.3\Big) = \frac{u-5}{20}+\frac{0.6}{u}. \nonumber
\end{align}
The first-order condition \eqref{r4e7-foc} is $D+N\big(1+2b'\big)=0$, i.e.\ $D+N-12N/u^2=0$:
\begin{align}
  \Big[\frac{u-5}{20}+\frac{0.6}{u}\Big]+\frac{5-u}{20}-\frac{12}{u^2}\cdot\frac{5-u}{20} &= 0 \nonumber\\
  \frac{0.6}{u}-\frac{12(5-u)}{20u^2} &= 0 \nonumber\\
  12u-12(5-u) &= 0 \qquad(\text{multiplying by }20u^2) \nonumber\\
  24u &= 60 \;\Rightarrow\; u=2.5,\quad c^*=\frac{2.5-1}{20}=0.075. \nonumber
\end{align}
At $c^*$: $b=0.3/2.5=0.12$; $D=0.4+0.075+2(0.12-0.3)=0.475-0.36=0.115$; $N=0.2-0.075=0.125$;
$k(c^*)=1/0.115=8.6957$; $(\rho_1-1-c^*)k(c^*)=0.125/0.115=1.0870$, against $0.5$ at $c=0$
(a grid search confirms $c^*$ is the global maximum). With $A=10$: $I=86.96$ and $U_b=10.87$
(instead of $25$ and $5$). Checks: $R_b=0.12/0.4=0.3$ and $R_m=0.075/0.4=0.1875$ per unit of $I$;
$I_m=(p_L/\Delta p)c^*I=(1)(0.075)(86.96)=6.52$; the uninformed investors pay
$86.96-10-6.52=70.43$ and expect $(0.8)(1.5-0.3-0.1875)(86.96)=70.43$.

\intuition Each unit of investment needs a stake $BI/\Delta p$ for the insider, so only $\rho_0$
per unit can be promised to outsiders; the borrower's cash fills the gap $1-\rho_0$ per unit,
hence the multiplier $1/(1-\rho_0)$. A monitor shrinks the insider's temptation, so her required
stake falls and pledgeable income rises; he must be given a stake too, but if he pays for it up
front, the borrower's only loss is the monitoring cost $c$ per unit. The best $c$ balances
cheaper incentives against that cost.
\pitfalls
\begin{itemize}
  \item Forgetting the monitor's cash contribution $I_m$: the multiplier then comes out as the
        statement's printed expression, which assumes the monitor keeps a rent.
  \item Keeping the borrower's stake at $BI/\Delta p$ after monitoring; under monitoring her
        incentive constraint uses $b(c)I$.
  \item Maximizing the multiplier $k(c)$ instead of $(\rho_1-1-c)k(c)$: the NPV per unit
        falls with $c$.
\end{itemize}
\bookans Matches the book's answer (p.~574): $(\Delta p)R_m\ge cI$, $(\Delta p)R_b\ge b(c)I$,
$I_m=p_HR_m-cI$, $k(c)=1/\{1-\rho_0+c+(p_H/\Delta p)[b(c)-B]\}$, and the borrower maximizes
$(\rho_1-1-c)k(c)$. The objective printed in the statement (p.~147) differs from this answer
by the monitor's contribution; see Step~8.
\refbib{Tirole Ch.~3, \S3.4.1--3.4.2 ``The Continuous-Investment Model'', ``The Lenders' Credit Analysis'', pp.~127--128; Ch.~9, \S9.2.1 ``Benefit of Activism'' (monitoring capital not scarce), pp.~356--358.}

% =============================================================================
\exsec{8. Lender market power --- Tirole Exercise 3.13, p.~149}

\asked (i) Fixed investment, a single lender who makes a take-it-or-leave-it offer: (i-a) the
contract the lender chooses; (i-b) whether the project is financed for the same values of $A$ as
with competitive lenders; (i-c) the borrower's net utility as a function of $A$ over four
regions, and why it is not monotonic. (ii) Variable investment: the lender's contract and its
effect on the scale of investment; show that investment falls.

\approach Same Holmstr\"om--Tirole moral hazard, but the roles are reversed: the lender
maximizes his profit, subject to the borrower's incentive constraint (IC$_b$) and her
participation constraint (IR$_b$: net utility $\ge 0$, since she can consume $A$). With
competitive lenders the investors' constraint binds and the borrower gets the NPV; here the
lender tries to take the NPV and is stopped only by the incentive rent he must leave her.

\notation
\begin{itemize}
  \item $I$: investment ($I>A$ in (i)); $A$: borrower's cash; $R$: income on success; $p_H$, $p_L$, $\Delta p$, $B$ as in Exercise~7.
  \item $V=p_HR-I>0$: NPV; $\bar A$: $p_H(R-B/\Delta p)=I-\bar A$; $\hat A=p_HB/\Delta p$ (both from the statement), $\bar A>0$.
  \item $R_b$: borrower's reward on success; $U_b=p_HR_b-A$: her net utility; $U_l=p_H(R-R_b)-(I-A)$: lender's profit.
  \item (ii): income $RI$, private benefit $BI$, $\rho_1=p_HR>1>\rho_0=p_H(R-B/\Delta p)$.
\end{itemize}

\step{1}{(i) The lender's program} If $A\ge I$ the borrower does not need the lender and gets
$U_b=V$. For $A<I$, the lender chooses $R_b$:
\begin{align}
  \max_{R_b}\ & U_l=p_H(R-R_b)-(I-A) \nonumber\\
  \text{s.t. }\ & (\Delta p)R_b\ge B, \tag{IC$_b$}\\
                & p_HR_b-A\ge 0. \tag{IR$_b$}
\end{align}

\step{2}{(i-a) The optimal reward} $U_l$ falls with $R_b$, so the lender picks the smallest
$R_b$ meeting both constraints: (IC$_b$) requires $R_b\ge B/\Delta p$ and (IR$_b$) requires
$R_b\ge A/p_H$. Hence
\begin{align}
  R_b^* &= \max\Big\{\frac{B}{\Delta p},\,\frac{A}{p_H}\Big\}, \label{r4e8-Rb}
\end{align}
and $A/p_H\ge B/\Delta p$ if and only if $A\ge p_HB/\Delta p=\hat A$.

\step{3}{(i-a) The lender's profit} For $A<\hat A$, $R_b^*=B/\Delta p$ and, using the definition
of $\bar A$,
\begin{align}
  U_l &= p_H\Big(R-\frac{B}{\Delta p}\Big)-I+A \nonumber\\
      &= (I-\bar A)-I+A = A-\bar A. \label{r4e8-Ul1}
\end{align}
For $A\ge\hat A$, $R_b^*=A/p_H$ and
\begin{align}
  U_l &= p_HR-p_H\frac{A}{p_H}-I+A \nonumber\\
      &= p_HR-I = V. \label{r4e8-Ul2}
\end{align}

\answer{(i-a)} The lender offers $R_b=\max\{B/\Delta p,\,A/p_H\}$ on success and $0$ on failure
(and lends only if \eqref{r4e8-Ul1} is nonnegative).

\step{4}{(i-b) The financing decision} The lender lends if $U_l\ge 0$. By \eqref{r4e8-Ul1}, for
$A<\hat A$ this holds if and only if $A\ge\bar A$; by \eqref{r4e8-Ul2}, for $A\ge\hat A$ it
always holds. The two regions fit together because
\begin{align}
  \hat A-\bar A &= \frac{p_HB}{\Delta p}-\Big[I-p_H\Big(R-\frac{B}{\Delta p}\Big)\Big]
                 = p_HR-I = V > 0. \nonumber
\end{align}
So the project is financed if and only if $A\ge\bar A$, exactly the competitive threshold (3.3).

\answer{(i-b)} No: the project is funded if and only if $A\ge\bar A$, as with competitive
lenders. Market power changes only who gets the NPV.

\step{5}{(i-c) The borrower's net utility} Substituting \eqref{r4e8-Rb} into $U_b=p_HR_b-A$
region by region (assuming $\hat A<I$, i.e.\ $p_HB/\Delta p<I$, so that the four regions exist):
\begin{align}
  U_b(A) &= \begin{cases}
     0 & A<\bar A \quad(\text{no loan}),\\[2pt]
     p_H\dfrac{B}{\Delta p}-A=\hat A-A & \bar A\le A<\hat A,\\[6pt]
     p_H\dfrac{A}{p_H}-A=0 & \hat A\le A<I,\\[2pt]
     V & A\ge I \quad(\text{self-financing}).
  \end{cases} \label{r4e8-Ub}
\end{align}
At $A=\bar A$, $U_b=\hat A-\bar A=V$: utility jumps from $0$ to the full NPV, then falls one for
one with $A$ to $0$ at $\hat A$, stays at $0$, and jumps back to $V$ at $A=I$
(Figure~\ref{r4e8-fig}).

\begin{figure}[ht]
\centering
\begin{tikzpicture}[x=0.075cm,y=0.12cm,>=stealth]
  \draw[->] (0,0) -- (125,0) node[right] {$A$};
  \draw[->] (0,0) -- (0,30) node[above] {$U_b$};
  \draw[very thick,blue!60!black] (0,0) -- (40,0);
  \draw[very thick,blue!60!black] (40,20) -- (60,0) -- (100,0);
  \draw[very thick,blue!60!black] (100,20) -- (122,20);
  \draw[dashed,gray] (40,0) -- (40,20) (100,0) -- (100,20);
  \draw[dotted,red!70!black,thick] (40,20) -- (122,20);
  \fill[blue!60!black] (40,20) circle (0.8pt) (100,20) circle (0.8pt);
  \draw[blue!60!black] (40,0) circle (0.8pt) (100,0) circle (0.8pt);
  \node[below] at (40,0) {$\bar A=40$};
  \node[below] at (60,0) {$\hat A=60$};
  \node[below] at (100,0) {$I=100$};
  \node[left] at (0,20) {$V=20$};
  \node[red!70!black,above] at (78,20) {\footnotesize competitive lenders};
\end{tikzpicture}
\caption{Borrower's net utility with a monopoly lender (solid) and with competitive lenders
(dotted), Exercise~8 illustration: $I=100$, $R=150$, $p_H=0.8$, $p_L=0.4$, $B=30$.}
\label{r4e8-fig}
\end{figure}

\answer{(i-c)} $U_b=0$ on $(-\infty,\bar A)$; $\hat A-A$ (decreasing from $V$ to $0$) on
$[\bar A,\hat A)$; $0$ on $[\hat A,I)$; $V$ on $[I,\infty)$. It is nonmonotonic because the
lender extracts everything above the incentive rent $p_HB/\Delta p$, and a richer borrower
(who must contribute her cash) has less of that rent left; only when she needs no loan does she
escape the lender.

\emph{Numerical illustration.} $I=100$, $R=150$, $p_H=0.8$, $p_L=0.4$, $B=30$ (assumption (3.2):
$0.4(150)+30-100=-10<0$). Then $V=0.8(150)-100=20$; $B/\Delta p=75$; $p_H(R-B/\Delta p)=0.8(75)=60$;
$\bar A=100-60=40$; $\hat A=0.8(75)=60$. At $A=50$: $R_b^*=\max\{75,\,50/0.8=62.5\}=75$,
$U_b=0.8(75)-50=10$, $U_l=0.8(150-75)-50=10$, total $20=V$ (with competitive lenders $U_b=20$).
At $A=70$: $R_b^*=\max\{75,87.5\}=87.5$, $U_b=0.8(87.5)-70=0$, $U_l=0.8(62.5)-30=50-30=20=V$.

\step{6}{(ii) The lender's program with variable investment} The lender now chooses the scale
$I$ and the reward $R_b$:
\begin{align}
  \max_{I,\,R_b}\ & U_l=p_H(RI-R_b)-(I-A) \nonumber\\
  \text{s.t. }\ & (\Delta p)R_b\ge BI, \tag{IC$_b$}\\
                & p_HR_b\ge A. \tag{IR$_b$}
\end{align}

\step{7}{(ii) Both constraints bind} The program is linear, so the Kuhn--Tucker conditions are
sufficient. The Lagrangian, with multipliers $\lambda,\mu\ge 0$, is
\begin{align}
  \mathcal L &= p_H(RI-R_b)-I+A+\lambda\big[(\Delta p)R_b-BI\big]+\mu\big[p_HR_b-A\big]. \nonumber
\end{align}
Setting the derivatives to zero,
\begin{align}
  \frac{\partial\mathcal L}{\partial I} &= p_HR-1-\lambda B=0
     \;\Rightarrow\; \lambda=\frac{\rho_1-1}{B}>0, \nonumber\\
  \frac{\partial\mathcal L}{\partial R_b} &= -p_H+\lambda\Delta p+\mu p_H=0
     \;\Rightarrow\; \mu=1-\frac{\lambda\Delta p}{p_H}=1-\frac{(\rho_1-1)\Delta p}{p_HB}. \nonumber
\end{align}
$\lambda>0$ because $\rho_1>1$, and
\begin{align}
  \mu>0 \iff \frac{(\rho_1-1)\Delta p}{p_HB}<1 \iff \rho_1-1<\frac{p_HB}{\Delta p}
        \iff \rho_1-\frac{p_HB}{\Delta p}<1 \iff \rho_0<1, \nonumber
\end{align}
which holds. Both multipliers are positive, so both constraints bind. (The book's verbal version:
if (IC$_b$) were slack, (IR$_b$) would bind and $U_l=(\rho_1-1)I$ would push $I$ to infinity,
violating (IC$_b$); if (IR$_b$) were slack, (IC$_b$) would bind and $U_l=(\rho_0-1)I+A$ would
push $I$ to $0$, violating (IR$_b$).)

\step{8}{(ii) Scale and profit} From the binding (IR$_b$), $R_b=A/p_H$; substituting into the
binding (IC$_b$),
\begin{align}
  (\Delta p)\frac{A}{p_H} &= BI \nonumber\\
  I^m &= \frac{\Delta p}{p_HB}A = \frac{A}{p_HB/\Delta p}. \label{r4e8-Im}
\end{align}
The lender's profit is $U_l=p_HRI^m-p_H(A/p_H)-I^m+A=(\rho_1-1)I^m$, and $U_b=p_HR_b-A=0$.

\step{9}{(ii) Comparison with competitive lenders} From Exercise~7(i), $I^*=A/(1-\rho_0)$. Then
\begin{align}
  I^m<I^* &\iff \frac{p_HB}{\Delta p}>1-\rho_0 = 1-\rho_1+\frac{p_HB}{\Delta p} \iff \rho_1>1, \nonumber
\end{align}
which holds. With the illustration of Exercise~7 ($p_H=0.8$, $\Delta p=0.4$, $R=1.5$, $B=0.3$,
$A=10$): $p_HB/\Delta p=0.8(0.3)/0.4=0.6$, so $I^m=10/0.6=16.667<I^*=25$; the multipliers are
$\lambda=0.2/0.3=0.6667$ and $\mu=1-0.2(0.4)/(0.8\cdot0.3)=0.6667$; $U_l=0.2(16.667)=3.3333$,
$U_b=0$. (The formula presumes $p_HB/\Delta p<1$, so that $I^m>A$; otherwise the borrower could
do better by investing her cash on her own, an outside option the program's (IR$_b$) ignores.)

\answer{(ii)} The lender sets $R_b=A/p_H=BI/\Delta p$ and $I^m=A\,\Delta p/(p_HB)$, keeping the
whole NPV $(\rho_1-1)I^m$. Since $I^m<A/(1-\rho_0)$, lender market power reduces the scale of
investment.

\intuition In both versions the lender would like to take everything, but he must leave the
borrower the stake $B/\Delta p$ (or $BI/\Delta p$) that keeps her honest. With a fixed project
that rent does not depend on any choice, so the financing decision is unchanged and only the
split moves; a richer borrower brings more cash, and the lender simply takes it. With a variable
project the rent grows with scale: each extra unit of investment forces the lender to leave
$p_HB/\Delta p$ more to the borrower. Like a monopolist restricting output, the lender stops
where the borrower's rent is exactly what she put in, below the competitive scale.
\pitfalls
\begin{itemize}
  \item Imposing the investors' breakeven constraint (competitive case) instead of the
        borrower's participation constraint: the lender's profit is not zero here.
  \item Concluding that market power must raise the financing threshold in the fixed model; it
        does not, because the lender gains from any loan with $A\ge\bar A$.
  \item Forgetting the region $A\ge I$, where the borrower self-finances and gets $V$.
\end{itemize}
\bookans Matches the book's answer (pp.~576--577): $R_b=A/p_H$ when this satisfies (IC$_b$),
i.e.\ $A\ge\hat A$, otherwise $R_b=B/\Delta p$ with $U_b=p_HB/\Delta p-A$; financing if and only if
$A\ge\bar A$; the four-region figure; in (ii) both constraints bind and
$I=A/(p_HB/\Delta p)<I^*=A/[1-p_H(R-B/\Delta p)]$.
\refbib{Tirole Ch.~3, \S3.2.1--3.2.2, pp.~115--117; \S3.4.1--3.4.2, pp.~127--128. Lecture Topic~4A.}

%% BEGIN-FROM lista-topic4-moral-hazard-overhang-solutions.tex EX 9 AS 9
\exsec{9. Entrepreneurial risk aversion --- Tirole Exercise 3.2, p.~144}

\asked (i) The entrepreneur is risk neutral above a subsistence level $c_0$ and infinitely averse
to falling below it. Find the smallest cash $\bar A$ with which risk-neutral investors finance the
project, and explain why the answer differs between a project that succeeds for sure under effort
($p_H=1$) and one that can fail even under effort ($p_H<1$). (ii) Redo the analysis for a general
strictly concave utility, once with limited liability (her payments cannot be negative) and once
without.

\approach This is the fixed-investment model of the lecture (Topic~4A) with a risk-averse
borrower. The contract pays her $c^S$ on success and $c^F$ on failure; she chooses it to maximise
expected utility subject to the investors' breakeven constraint (IR$_l$) and her incentive
constraint (IC$_b$), now written in utility. Pledgeable income is total expected income minus the
cheapest expected payment that satisfies (IC$_b$) and the lower bound on $c^F$; a floor on
consumption (subsistence in (i), limited liability in (ii)) raises that payment and so the minimum
equity. In (ii) we write the Lagrangian and derive the first-order conditions, which give the
Holmstr\"om--Shavell insurance-versus-incentives trade-off.

\notation
\begin{itemize}
  \item $A$: entrepreneur's cash ($A\ge c_0$ in (i)); $I>A$: fixed investment; $R>0$: income in
        the case of success, $0$ in the case of failure.
  \item $p_H$, $p_L=p_H-\Delta p$: success probabilities when she works / shirks ($0<p_L<p_H\le 1$);
        $B>0$: private benefit from shirking.
  \item $\mathit{NPV}\equiv p_HR-I>0$; $\bar A_{rn}\equiv I-p_H(R-B/\Delta p)>0$: the
        risk-neutral threshold (3.3), positive by the assumption $I>p_H(R-B/\Delta p)$.
  \item $c^S$, $c^F$: the entrepreneur's consumption in the case of success / failure (her total
        payoff from the contract); $u(\cdot)$: utility of consumption; $c_0$: subsistence level in (i).
  \item $\mathcal{P}$: expected pledgeable income; $\bar A$: minimum cash for financing.
  \item $h(c)\equiv u^{-1}\!\big(u(c)+B/\Delta p\big)$: consumption in success that gives exactly the
        utility premium $B/\Delta p$ over consumption $c$ in failure (part (ii)).
\end{itemize}
\assume The entrepreneur puts all of $A$ into the project and the contract pays her $c^S$ or $c^F$
(keeping cash aside is the same as adding it to both payments). In (ii) the private benefit enters
utility additively, $p\,u(c^S)+(1-p)\,u(c^F)+B$ when she shirks, as in the book's answer.

\step{1}{(i), $p_H<1$: the programme} The entrepreneur offers the contract that maximises her
expected utility subject to the investors' breakeven and her incentive constraint:
\begin{align}
  \max_{c^S,c^F}\; & p_H\,u(c^S)+(1-p_H)\,u(c^F) \label{t4e9-obj}\\
  \text{s.t. }\; & p_H(R-c^S)+(1-p_H)(0-c^F)\ \ge\ I-A, \label{t4e9-irl}\\
  & p_H\,u(c^S)+(1-p_H)\,u(c^F)\ \ge\ p_L\,u(c^S)+(1-p_L)\,u(c^F)+B. \label{t4e9-ic}
\end{align}
Failure has probability $1-p_H>0$ even under work, so $c^F<c_0$ would make \eqref{t4e9-obj} equal to
$-\infty$: the optimum has $c^F\ge c_0$, and then $c^S\ge c_0$ too. On $[c_0,\infty)$ utility is
linear, $u(c)=c$, so \eqref{t4e9-ic} reads
\begin{align}
  p_Hc^S+(1-p_H)c^F &\ge p_Lc^S+(1-p_L)c^F+B \nonumber\\
  (p_H-p_L)c^S-\big[(1-p_L)-(1-p_H)\big]c^F &\ge B \nonumber\\
  (p_H-p_L)c^S-(p_H-p_L)c^F &\ge B \nonumber\\
  (\Delta p)(c^S-c^F) &\ge B. \label{t4e9-iclin}
\end{align}

\step{2}{(i), $p_H<1$: which constraints bind, and $\bar A$} With linear utility the objective
\eqref{t4e9-obj} is the expected payment $p_Hc^S+(1-p_H)c^F$. If \eqref{t4e9-irl} were slack, raising
$c^S$ would raise the objective and relax \eqref{t4e9-iclin}; so \eqref{t4e9-irl} binds and the
entrepreneur gets $A+\mathit{NPV}$. Financing is possible if and only if the cheapest contract that
satisfies \eqref{t4e9-iclin} and $c^F\ge c_0$ leaves the investors enough. Using
$c^S\ge c^F+B/\Delta p$ from \eqref{t4e9-iclin} and then $c^F\ge c_0$, any feasible contract costs
\begin{align}
  p_Hc^S+(1-p_H)c^F &\ge p_H\Big(c^F+\frac{B}{\Delta p}\Big)+(1-p_H)c^F \nonumber\\
                    &= c^F+p_H\frac{B}{\Delta p} \ \ge\ c_0+p_H\frac{B}{\Delta p}, \nonumber
\end{align}
with equality iff both inequalities bind, so the cheapest contract is
\begin{align}
  c^F &= c_0, \label{t4e9-cf}\\
  c^S &= c_0+\frac{B}{\Delta p}. \label{t4e9-cs}
\end{align}
The pledgeable income is total expected income minus this minimum expected payment; substituting
\eqref{t4e9-cf}--\eqref{t4e9-cs},
\begin{align}
  \mathcal{P} &= p_HR-\Big[p_H\Big(c_0+\frac{B}{\Delta p}\Big)+(1-p_H)(c_0)\Big] \nonumber\\
              &= p_H\Big(R-\frac{B}{\Delta p}\Big)-c_0. \label{t4e9-P}
\end{align}
(the middle step: $p_Hc_0+(1-p_H)c_0=c_0$). Financing requires $\mathcal{P}\ge I-A$; using
\eqref{t4e9-P} and moving terms,
\begin{align}
  p_H\Big(R-\frac{B}{\Delta p}\Big)-c_0 &\ge I-A \nonumber\\
  A \ \ge\ \bar A &= I+c_0-p_H\Big(R-\frac{B}{\Delta p}\Big) \label{t4e9-Abar}\\
                  &= \Big[I-p_H\Big(R-\frac{B}{\Delta p}\Big)\Big]+c_0=\bar A_{rn}+c_0. \nonumber
\end{align}
Because $\bar A_{rn}>0$, $\bar A>c_0$: entrepreneurs with $A\in[c_0,\bar A)$ are rationed. The
participation constraint holds, since she gets $A+\mathit{NPV}>A$. Illustration: with
$p_H=0.8$, $p_L=0.4$, $B=1$, $R=10$, $I=6.5$, $c_0=1$, \eqref{t4e9-Abar} gives
$\bar A=(6.5)+(1)-(0.8)(10-1/0.4)=1.5$, against $\bar A_{rn}=0.5$.

\step{3}{(i), $p_H=1$} Now failure never happens on the equilibrium path; it happens only if she
shirks, with probability $1-p_L>0$. Set $c^F<c_0$ (for example $c^F=0$, feasible under limited
liability when $c_0>0$). Shirking then yields $p_L\,u(c^S)+(1-p_L)(-\infty)+B=-\infty$, so
\eqref{t4e9-ic} holds for every $c^S\ge c_0$ and the incentive problem costs nothing. The cheapest
contract is $c^S=c_0$, and
\begin{align}
  \mathcal{P} &= R-c_0, \label{t4e9-P1}\\
  \bar A &= I+c_0-R\ <\ c_0, \label{t4e9-Abar1}
\end{align}
where the inequality uses $R-I=\mathit{NPV}>0$. Every entrepreneur with $A\ge c_0$ is financed: no
rationing (with the numbers above, $\bar A=6.5+1-10=-2.5$). The threshold is discontinuous at
$p_H=1$: by \eqref{t4e9-Abar}, $\bar A\to I+c_0-R+B/\Delta p$ as $p_H\to1$, a jump of $B/\Delta p$.
If $c_0\le0$, limited liability rules out a punishment below $c_0$ and the risk-neutral analysis applies.

\answer{(i)} For $p_H<1$: $\bar A=I+c_0-p_H(R-B/\Delta p)=\bar A_{rn}+c_0$, so entrepreneurs with
$A\in[c_0,\bar A)$ are rationed. For $p_H=1$: failure reveals shirking and can be punished with
$c^F<c_0$ at no cost on the equilibrium path, so $\bar A=I+c_0-R<c_0$ and every entrepreneur with
$A\ge c_0$ is financed (no rationing).

\step{4}{(ii) The general programme} With $u'>0$, $u''<0$, the contract solves
\begin{align}
  \max_{c^S,c^F}\; & p_H\,u(c^S)+(1-p_H)\,u(c^F) \label{t4e9-gobj}\\
  \text{s.t. }\; & p_H(R-c^S)-(1-p_H)c^F\ \ge\ I-A, \tag{IR$_l$}\label{t4e9-girl}\\
  & (\Delta p)\big[u(c^S)-u(c^F)\big]\ \ge\ B, \tag{IC$_b$}\label{t4e9-gic}\\
  & c^F\ge 0,\ c^S\ge0 \quad\text{(only under limited liability)}, \tag{LL}\label{t4e9-ll}\\
  & p_H\,u(c^S)+(1-p_H)\,u(c^F)\ \ge\ u(A). \tag{PC}\label{t4e9-pc}
\end{align}
\eqref{t4e9-gic} is the risk-averse version of (IC$_b$): subtract $p_Lu(c^S)+(1-p_L)u(c^F)+B$
from \eqref{t4e9-gobj}. \eqref{t4e9-pc} says she prefers the project to consuming $A$. In utility
units $v^S=u(c^S)$, $v^F=u(c^F)$ the objective and \eqref{t4e9-gic} are linear and the cost
$p_Hu^{-1}(v^S)+(1-p_H)u^{-1}(v^F)$ is convex, so the problem is convex and the first-order
conditions are sufficient.

\step{5}{(ii) Which constraints bind} We show both (IR$_l$) and (IC$_b$) bind, the second by
deriving the first-order conditions. (a) \eqref{t4e9-girl} binds: otherwise raise $c^S$, which raises
\eqref{t4e9-gobj} and relaxes \eqref{t4e9-gic}. (b) With multipliers
$\lambda\ge0$ on \eqref{t4e9-girl} and $\mu\ge0$ on \eqref{t4e9-gic} (and \eqref{t4e9-ll} slack), the
Lagrangian is
\begin{align}
  \mathcal L &= p_H\,u(c^S)+(1-p_H)\,u(c^F)+\lambda\big[p_H(R-c^S)-(1-p_H)c^F-I+A\big] \nonumber\\
             &\quad+\mu\big[(\Delta p)\big(u(c^S)-u(c^F)\big)-B\big]. \nonumber
\end{align}
Differentiating with respect to $c^S$ and $c^F$ and setting the derivatives to zero,
\begin{align}
  \frac{\partial\mathcal L}{\partial c^S} &= p_Hu'(c^S)-\lambda p_H+\mu(\Delta p)u'(c^S)=0, \nonumber\\
  &\Longrightarrow\quad u'(c^S)\big(p_H+\mu\Delta p\big)=\lambda p_H, \nonumber\\
  \frac{\partial\mathcal L}{\partial c^F} &= (1-p_H)u'(c^F)-\lambda(1-p_H)-\mu(\Delta p)u'(c^F)=0, \nonumber\\
  &\Longrightarrow\quad u'(c^F)\big(1-p_H-\mu\Delta p\big)=\lambda(1-p_H). \nonumber
\end{align}
Here $\lambda>0$, because $\lambda=0$ would make the first condition read $u'(c^S)(p_H+\mu\Delta p)=0$,
impossible with $u'>0$. Dividing the first by $\lambda p_H\,u'(c^S)$ and the second by
$\lambda(1-p_H)\,u'(c^F)$,
\begin{align}
  \frac{1}{u'(c^S)} &= \frac{1}{\lambda}\Big(1+\mu\frac{\Delta p}{p_H}\Big), \label{t4e9-focs}\\
  \frac{1}{u'(c^F)} &= \frac{1}{\lambda}\Big(1-\mu\frac{\Delta p}{1-p_H}\Big). \label{t4e9-focf}
\end{align}
If $\mu=0$, \eqref{t4e9-focs}--\eqref{t4e9-focf} give $u'(c^S)=u'(c^F)$, i.e.\ full insurance
$c^S=c^F$, which violates \eqref{t4e9-gic} because $B>0$. So $\mu>0$ and, by \eqref{t4e9-focs}--\eqref{t4e9-focf},
$c^S>c^F$: the entrepreneur bears risk, but only as much as incentives require
(Holmstr\"om--Shavell inverse-marginal-utility rule).

\step{6}{(ii) The solution} We reduce the two binding constraints to one equation in $c^F$. The
binding \eqref{t4e9-gic} is $u(c^S)=u(c^F)+B/\Delta p$, i.e.\ $c^S=h(c^F)$. The binding
\eqref{t4e9-girl} is $p_HR-p_Hc^S-(1-p_H)c^F=I-A$, i.e.
\begin{align}
  p_Hc^S+(1-p_H)c^F=A+(p_HR-I)=A+\mathit{NPV}. \nonumber
\end{align}
Substituting $c^S=h(c^F)$,
\begin{align}
  p_H\,h(c^F)+(1-p_H)\,c^F = A+\mathit{NPV}. \label{t4e9-sol}
\end{align}
Differentiating the identity $u(h(c))=u(c)+B/\Delta p$ with respect to $c$,
\begin{align}
  u'\big(h(c)\big)\,h'(c)=u'(c) \quad\Longrightarrow\quad h'(c)=\frac{u'(c)}{u'(h(c))}\ \ge\ 1, \nonumber
\end{align}
because $h(c)>c$ and $u'$ is decreasing. Hence the left-hand side of \eqref{t4e9-sol} is increasing
in $c^F$, so \eqref{t4e9-sol} defines a unique $c^F=c(A)$, increasing in $A$. Substituting $u(c^S)=u(c^F)+B/\Delta p$ into
\eqref{t4e9-gobj}, her utility is
\begin{align}
  V(A) &= p_H\Big(u(c(A))+\frac{B}{\Delta p}\Big)+(1-p_H)\,u(c(A)) \nonumber\\
       &= u(c(A))+\frac{p_HB}{\Delta p}. \label{t4e9-V}
\end{align}

\step{7}{(ii) Without limited liability} Only \eqref{t4e9-pc} can fail. Let $\mathit{CE}$ be the certainty
equivalent of the lottery $(c^S,c^F)$, $u(\mathit{CE})=V(A)$, and $\pi\equiv \mathbb{E}[c]-\mathit{CE}$ its
risk premium; by \eqref{t4e9-sol}, $\mathbb{E}[c]=A+\mathit{NPV}$. Then
\begin{align}
  \text{financing}\iff V(A)\ge u(A)\iff \mathit{CE}\ge A\iff \pi(A)\le \mathit{NPV}. \label{t4e9-nll}
\end{align}
A risk-neutral entrepreneur ($\pi=0$) without limited liability is always financed (she can be
punished with $c^F<0$; there is no $\bar A$). A risk-averse one is financed only if the risk premium
of the incentive lottery does not exceed the NPV: risk aversion alone can block positive-NPV projects.
Whether this gives a threshold in $A$ depends on how $\pi$ varies with wealth (with $B$ in utility
units, a richer entrepreneur needs a larger money spread to produce the utility gap $B/\Delta p$).

\step{8}{(ii) With limited liability} \eqref{t4e9-ll} requires $c(A)\ge0$. Because the left-hand side of
\eqref{t4e9-sol} is increasing, $c(A)\ge0$ iff $A+\mathit{NPV}\ge p_Hh(0)$, i.e.
\begin{align}
  A\ \ge\ A_{LL}\equiv I-p_H\big[R-h(0)\big]. \label{t4e9-all}
\end{align}
Below $A_{LL}$ no contract satisfies \eqref{t4e9-girl}, \eqref{t4e9-gic} and \eqref{t4e9-ll}: the cheapest
incentive-compatible contract, $c^F=0$, $c^S=h(0)$, already costs more than $A+\mathit{NPV}$. Above
$A_{LL}$ the solution is that of Step 6 with $c^F=c(A)\ge0$, so the minimum equity is
$\bar A=\min\{A\ge A_{LL}:\ V(A)\ge u(A)\}$. Two remarks. First, for $A>A_{LL}$ the entrepreneur
receives $c^F>0$ in failure: unlike the risk-neutral case, pure debt ($c^F=0$) is no longer optimal,
because she buys partial insurance. Second, risk neutrality gives back the book: with $u(c)=c$,
$h(0)=B/\Delta p$ and \eqref{t4e9-all} is \eqref{t4e9-Abar} with $c_0=0$; part (i) is the case of a floor
$c_0$ instead of $0$. If $u$ is normalised so that $u'(0)=1$, concavity gives $h(0)\ge B/\Delta p$, so
$A_{LL}\ge\bar A_{rn}$: risk aversion raises the equity requirement.

\answer{(ii)} Both (IR$_l$) and (IC$_b$) bind; the contract solves \eqref{t4e9-sol} with
$u(c^S)=u(c^F)+B/\Delta p$, and the first-order conditions \eqref{t4e9-focs}--\eqref{t4e9-focf} give
$c^S>c^F$ (partial insurance). Without limited liability, financing obtains iff the risk premium of
the incentive lottery is at most the NPV, $\pi(A)\le\mathit{NPV}$. With limited liability, also
$A\ge A_{LL}=I-p_H[R-h(0)]$, and $\bar A=\min\{A\ge A_{LL}:V(A)\ge u(A)\}$; above $A_{LL}$ she gets
$c^F>0$, so pure debt is no longer optimal.

\intuition Incentives need a spread between the success and failure payoffs. A risk-neutral
entrepreneur bears that spread for free; a subsistence floor or limited liability bounds the
punishment, so the spread must be paid as a rent in success, and a risk-averse entrepreneur must in
addition be compensated for the risk. Both raise the stake she keeps and lower pledgeable income.
The $p_H=1$ case is the exception: failure then proves shirking, so an unbounded punishment off the
equilibrium path removes the agency cost.

\pitfalls
\begin{itemize}
  \item Subtracting $c_0$ from the cash instead of adding it to the payments: the entrepreneur must be
        paid at least $c_0$ in \emph{both} states, which costs investors $c_0$ in expectation; the two
        bookkeepings agree only if this is done consistently.
  \item Treating $p_H=1$ like $p_H<1$. With $p_H=1$ failure never occurs under effort, so the
        $-\infty$ punishment is free; with $p_H<1$ it would hit her on the equilibrium path.
  \item In (ii), carrying over the risk-neutral result that pure debt ($c^F=0$) is optimal. A
        risk-averse borrower buys partial insurance, $c^F>0$, whenever $A>A_{LL}$.
\end{itemize}

\bookans Matches the book's answer (p.~571): $\bar A=I+c_0-p_H(R-B/\Delta p)$ for $p_H<1$, the
programme of (ii) with binding incentive constraint. For $p_H=1$ the book writes the pledgeable
income as $p_HR$, i.e.\ with $c_0$ kept aside from $A$; this is \eqref{t4e9-P1}--\eqref{t4e9-Abar1}
in other bookkeeping. The book only states the programme for (ii) and cites Holmstr\"om (1979) and
Shavell (1979); Steps 6--8 are this solution's.

\refbib Tirole \S3.2.1--3.2.2, pp.~115--117; answer p.~571.
%% END-FROM

%% BEGIN-FROM lista-topic4-moral-hazard-overhang-solutions.tex EX 10 AS 10
\exsec{10. Maximal incentives principle in the fixed-investment model --- Tirole Exercise 3.18, p.~151}

\asked In the fixed-investment model, failure now yields a positive income $R^F$ (salvage value)
instead of $0$. Show that a debt contract (investors take everything in failure and a fixed amount
in success) is an optimal financial structure, and show that, unlike the variable-investment model
of Section~3.4.3, other contracts may be optimal too.

\approach This is the lecture's ``debt dominates equity'' result (maximal incentives principle,
Topic~4A). Set up the borrower's programme: maximise her expected reward subject to (IC$_b$),
(IR$_l$) and limited liability. With a fixed investment her utility equals the NPV for \emph{every}
contract on which (IR$_l$) binds, so the contract matters only for feasibility, i.e.\ for
\emph{pledgeable income}. Paying her in failure buys no incentive, so the contract that maximises
pledgeable income pays her nothing in failure: debt. Uniqueness fails because, when $A>\bar A$, the
spare pledgeable income can be paid in either state.

\notation
\begin{itemize}
  \item $I$: fixed investment; $A$: entrepreneur's cash; $R^S>R^F>0$: income in success / failure,
        with $R^F<I-A$; $R\equiv R^S-R^F$.
  \item $p_H$, $p_L=p_H-\Delta p$, $B$: as in Section 3.2.
  \item $R_b^S$, $R_b^F\ge0$: entrepreneur's reward in success / failure (limited liability).
  \item $\mathit{NPV}\equiv p_HR^S+(1-p_H)R^F-I=p_HR+R^F-I>0$; $U_b$: entrepreneur's net utility.
  \item $D$: face value of debt; $\bar A$: minimum cash for financing.
\end{itemize}

\step{1}{Programme} The borrower designs the contract $(R_b^S,R_b^F)$; her incentive constraint is
$p_HR_b^S+(1-p_H)R_b^F\ge p_LR_b^S+(1-p_L)R_b^F+B$, which simplifies (collect $R_b^S$ and $R_b^F$
terms, $(1-p_L)-(1-p_H)=\Delta p$) to \eqref{t4e10-ic}:
\begin{align}
  U_b=\max_{R_b^S,R_b^F}\; & p_HR_b^S+(1-p_H)R_b^F-A \label{t4e10-obj}\\
  \text{s.t. }\; & (\Delta p)(R_b^S-R_b^F)\ \ge\ B, \tag{IC$_b$}\label{t4e10-ic}\\
  & p_H(R^S-R_b^S)+(1-p_H)(R^F-R_b^F)\ \ge\ I-A, \tag{IR$_l$}\label{t4e10-irl}\\
  & R_b^S\ge0,\ R_b^F\ge0. \tag{LL}\label{t4e10-ll}
\end{align}

\step{2}{The breakeven constraint binds; $U_b$ is the NPV} If \eqref{t4e10-irl} were slack, raising
$R_b^S$ would raise \eqref{t4e10-obj} and relax \eqref{t4e10-ic}. With \eqref{t4e10-irl} binding,
solve it for her expected reward,
\begin{align}
  p_HR^S-p_HR_b^S+(1-p_H)R^F-(1-p_H)R_b^F &= I-A \nonumber\\
  p_HR_b^S+(1-p_H)R_b^F &= p_HR^S+(1-p_H)R^F-(I-A), \nonumber
\end{align}
and substitute it into \eqref{t4e10-obj}:
\begin{align}
  U_b &= \big[p_HR^S+(1-p_H)R^F-(I-A)\big]-A \nonumber\\
      &= p_HR^S+(1-p_H)R^F-I \nonumber\\
      &= \mathit{NPV}. \label{t4e10-Ub}
\end{align}
Because $I$ is fixed, \eqref{t4e10-Ub} is the same for \emph{every} feasible contract with binding
\eqref{t4e10-irl}. The financial structure matters only for feasibility: it must deliver
\eqref{t4e10-ic} while leaving investors $I-A$.

\step{3}{Maximal pledgeable income and the financing condition} The constraint set is nonempty iff
the largest expected income that can go to investors subject to \eqref{t4e10-ic} and
\eqref{t4e10-ll} covers $I-A$, i.e.\ iff expected income minus the \emph{smallest} expected reward
compatible with \eqref{t4e10-ic} and \eqref{t4e10-ll} is at least $I-A$. Using
$R_b^S\ge R_b^F+B/\Delta p$ from \eqref{t4e10-ic},
\begin{align}
  p_HR_b^S+(1-p_H)R_b^F &\ge p_H\Big(R_b^F+\frac{B}{\Delta p}\Big)+(1-p_H)R_b^F \nonumber\\
                        &= R_b^F+p_H\frac{B}{\Delta p}\ \ge\ p_H\frac{B}{\Delta p}, \nonumber
\end{align}
the last step by $R_b^F\ge0$; both bounds are attained by $R_b^F=0$, $R_b^S=B/\Delta p$. Each unit of
failure reward costs one unit of expected income and buys no incentive. So
\begin{align}
  \mathcal{P} &= p_HR^S+(1-p_H)R^F-p_H\frac{B}{\Delta p} \nonumber\\
              &= p_H(R^S-R^F)+R^F-p_H\frac{B}{\Delta p} \nonumber\\
              &= R^F+p_H\Big(R-\frac{B}{\Delta p}\Big), \label{t4e10-P}
\end{align}
and financing obtains iff $\mathcal{P}\ge I-A$, i.e.
\begin{align}
  A\ \ge\ \bar A\equiv I-R^F-p_H\Big(R-\frac{B}{\Delta p}\Big)=p_H\frac{B}{\Delta p}-\mathit{NPV}, \label{t4e10-Abar}
\end{align}
where the second form uses $\mathit{NPV}=p_HR+R^F-I$.
The safe income $R^F$ works exactly like cash on hand.

\step{4}{A debt contract is optimal} Let investors hold debt of face value $D$ defined by
\begin{align}
  p_HD+(1-p_H)R^F = I-A. \label{t4e10-D}
\end{align}
In failure the firm defaults and investors take $R^F$; in success they get $D$. By $R^F<I-A$,
\eqref{t4e10-D} gives $D>R^F$, so the borrower gets $R_b^F=0$. Substituting \eqref{t4e10-D} into
\eqref{t4e10-Abar} (written as $p_H(R^S-B/\Delta p)+(1-p_H)R^F\ge I-A$),
\begin{align}
  p_H\Big(R^S-\frac{B}{\Delta p}\Big)+(1-p_H)R^F &\ge p_HD+(1-p_H)R^F \nonumber\\
  p_H\Big(R^S-\frac{B}{\Delta p}\Big) &\ge p_HD \nonumber\\
  R^S-\frac{B}{\Delta p} &\ge D \nonumber\\
  R_b^S=R^S-D &\ge \frac{B}{\Delta p}, \label{t4e10-debtic}
\end{align}
(subtracting $(1-p_H)R^F$, then dividing by $p_H>0$, then rearranging)
so \eqref{t4e10-ic} holds, \eqref{t4e10-irl} binds by construction, and by \eqref{t4e10-Ub} debt
attains the maximum. It is optimal whenever financing is feasible.

\step{5}{Debt is not uniquely optimal} The optimal set is every $(R_b^S,R_b^F)\ge0$ with binding
\eqref{t4e10-irl} and \eqref{t4e10-ic}. From the binding \eqref{t4e10-irl},
$R_b^S=[A+\mathit{NPV}-(1-p_H)R_b^F]/p_H$. Substituting this into \eqref{t4e10-ic},
\begin{align}
  (\Delta p)\Big(\frac{A+\mathit{NPV}-(1-p_H)R_b^F}{p_H}-R_b^F\Big) &\ge B \nonumber\\
  (\Delta p)\,\frac{A+\mathit{NPV}-(1-p_H)R_b^F-p_HR_b^F}{p_H} &\ge B \nonumber\\
  \frac{A+\mathit{NPV}-R_b^F}{p_H} &\ge \frac{B}{\Delta p} \nonumber\\
  A+\mathit{NPV}-R_b^F &\ge p_H\frac{B}{\Delta p} \nonumber\\
  R_b^F &\le A+\mathit{NPV}-p_H\frac{B}{\Delta p}=A-\bar A, \label{t4e10-set}
\end{align}
using \eqref{t4e10-Abar} in the last line. So every $R_b^F\in[0,A-\bar A]$ is optimal (add
$R_b^F\le R^F$ if investors' claim must be nonnegative in each state). The contract is unique only
at $A=\bar A$. For $A>\bar A$, a small failure reward (debt plus a little equity, or even all-equity
when its implied $R_b^F$ satisfies \eqref{t4e10-set}) does as well. Illustration: $R^S=12$, $R^F=2$,
$p_H=0.8$, $p_L=0.4$, $B=1.6$, $I=9$, $A=3$. Then $B/\Delta p=1.6/0.4=4$,
$\mathit{NPV}=0.8(12)+0.2(2)-9=9.6+0.4-9=1$, $\bar A=0.8(4)-1=2.2$, and by \eqref{t4e10-D}
$D=(9-3-0.2(2))/0.8=5.6/0.8=7$ (with $2<7\le12-4=8$). The optimal set is
$R_b^F\in[0,\,3-2.2]=[0,0.8]$; at $R_b^F=0.8$, $R_b^S=(3+1-0.2(0.8))/0.8=3.84/0.8=4.8$ and
\eqref{t4e10-ic} binds: $0.4(4.8-0.8)=1.6=B$.

\step{6}{Contrast with variable investment} In Section 3.4.3 the entrepreneur's utility
$(p_HR+R^F-1)I$ rises with $I$, so she expands $I$ until the financing constraint binds, the
analogue of $A=\bar A$. Then the slack in \eqref{t4e10-set} is zero, $R_b^F=0$ is forced and debt is
the unique optimum; any $R_b^F>0$ would waste incentive power and reduce the borrowing capacity (3.16).

\answer{(3.18)} Debt with face value $D$ given by $p_HD+(1-p_H)R^F=I-A$ (so $R^F<D\le R^S-B/\Delta p$)
is optimal whenever the project can be financed, $A\ge\bar A=I-R^F-p_H(R-B/\Delta p)$. It is not
uniquely optimal: every contract with $R_b^F\in[0,A-\bar A]$ and binding (IR$_l$) gives the same
utility $\mathit{NPV}$; the optimum is unique only at $A=\bar A$.

\intuition Paying the entrepreneur in failure costs expected income but weakens incentives, so it is
the worst way to reward her; debt, which gives her nothing in failure, maximises incentives per unit
of pay. With a fixed project, once the project is funded there is no further use for pledgeable
income, so any leftover slack can be paid out in either state.

\pitfalls
\begin{itemize}
  \item Concluding that all-equity is never optimal, as in Section~3.4.3. That needs the investment to
        expand until the financing constraint binds; with a fixed $I$ and $A>\bar A$ it does not bind.
  \item Leaving $R^F$ out of pledgeable income. The salvage value is not subject to moral hazard and
        works exactly like cash on hand: $\bar A$ falls one for one with $R^F$.
  \item Writing (IC$_b$) as $(\Delta p)R_b^S\ge B$, forgetting that a failure reward $R_b^F$ weakens
        incentives: what matters is the spread $R_b^S-R_b^F$.
\end{itemize}

\bookans Matches the book's answer (p.~579), including condition (1) there,
$p_H(R-B/\Delta p)\ge I-(A+R^F)$, the debt implementation with $R^F<D<R^S$ and the
non-uniqueness when (IR$_l$) holds with slack at $R_b^F=0$; \eqref{t4e10-set} makes the size of the
optimal set explicit.

\refbib Tirole \S3.2.2, pp.~116--117; \S3.4.3, pp.~128--129; answer p.~579.
%% END-FROM

%% BEGIN-FROM lista-topic4-moral-hazard-overhang-solutions.tex EX 11 AS 11
\exsec{11. Retained-earnings benefit --- Tirole Exercise 3.11, p.~148}

\asked (i) Find how much one extra unit of retained earnings (the entrepreneur's date-1 profit kept
in the firm) is worth to her at date 2, given that it helps finance a second project whose private
benefit $B^2$ is still random. (ii) Show, by argument or example, that the existence of the second
project can get the first project financed even though no date-2 income can be pledged at date 1.

\approach Two applications of the fixed-investment model (lecture Topic~4A, investment and internal
funds). At date 2, retained earnings play the role of cash $A$ in (IR$_l$): more of them lets the
second project be financed for higher values of $B^2$. Their \emph{shadow value} is the derivative of
the entrepreneur's gross date-2 utility. At date 1, that continuation value is lost when project 1
fails, so it works as an extra reward for effort: the date-1 incentive constraint needs less cash
reward, and date-1 \emph{pledgeable income} rises.

\notation
\begin{itemize}
  \item Project $t\in\{1,2\}$: investment $I^t$, income $R^t$ in success, probabilities $p_H^t$,
        $p_L^t=p_H^t-\Delta p^t$, private benefit $B^t$; $\mathit{NPV}^2\equiv p_H^2R^2-I^2>0$.
  \item $A^1$: entrepreneur's cash at date 1; $R_b^1$: her date-1 reward in success (her retained
        earnings), $0$ in failure.
  \item $B^2\in[\underline B^2,\bar B^2]$: random, c.d.f.\ $F$, density $f$, revealed at the start of date 2.
  \item $\hat B^2(R)$: largest $B^2$ for which project 2 is financed with retained earnings $R$.
  \item $W(R)$: entrepreneur's expected gross date-2 utility with retained earnings $R$; $\mu(R)\equiv W'(R)$:
        shadow value of retained earnings ($v$ in \S3.4.2).
\end{itemize}
\assume No discounting; the entrepreneur's only date-2 wealth is $R_b^1$ (or $0$), and date-2 lenders
cannot be pledged date-1 income.

\step{1}{(i) Date-2 financing given retained earnings $R$} Project 2 is financed iff its pledgeable
income covers the outside funding, (IR$_l$) for project 2:
\begin{align}
  p_H^2\Big(R^2-\frac{B^2}{\Delta p^2}\Big) &\ge I^2-R. \label{t4e11-f2}
\end{align}
Solving \eqref{t4e11-f2} for $B^2$ (expand, move $p_H^2B^2/\Delta p^2$ to the right and $I^2-R$ to
the left, then multiply by $\Delta p^2/p_H^2>0$),
\begin{align}
  p_H^2R^2-\frac{p_H^2B^2}{\Delta p^2} &\ge I^2-R \nonumber\\
  \frac{p_H^2B^2}{\Delta p^2} &\le p_H^2R^2-I^2+R=\mathit{NPV}^2+R \nonumber\\
  B^2\ \le\ \hat B^2(R) &\equiv \frac{\Delta p^2}{p_H^2}\big(\mathit{NPV}^2+R\big). \label{t4e11-Bhat}
\end{align}
By the assumption $\underline B^2>\Delta p^2\,\mathit{NPV}^2/p_H^2$, \eqref{t4e11-Bhat} gives
$\hat B^2(0)<\underline B^2$: with no retained earnings, project 2 is never financed.

\step{2}{(i) Gross date-2 utility} We follow the hint and compute her expected gross utility. If
financed, project 2's lenders break even and the entrepreneur gets her cash plus $\mathit{NPV}^2$
(equation (3.6)); otherwise she keeps $R$. With probability $F(\hat B^2(R))$ it is financed, so
\begin{align}
  W(R) &= F\big(\hat B^2(R)\big)\big(R+\mathit{NPV}^2\big)+\Big[1-F\big(\hat B^2(R)\big)\Big]R \nonumber\\
       &= R+F\big(\hat B^2(R)\big)\,\mathit{NPV}^2. \label{t4e11-W}
\end{align}

\step{3}{(i) Shadow value} Differentiating \eqref{t4e11-W} and using
$d\hat B^2/dR=\Delta p^2/p_H^2$ from \eqref{t4e11-Bhat},
\begin{align}
  \mu(R)=W'(R) &= 1+f\big(\hat B^2(R)\big)\,\frac{d\hat B^2}{dR}\,\mathit{NPV}^2 \nonumber\\
               &= 1+f\big(\hat B^2(R)\big)\,\frac{\Delta p^2}{p_H^2}\,\mathit{NPV}^2\ \ge 1. \label{t4e11-mu}
\end{align}
One unit of retained earnings is worth one unit of consumption plus the expected NPV it unlocks: it
raises the probability of financing at the margin by $f\,\Delta p^2/p_H^2$.

\answer{(i)} $\mu(R)=1+f\big(\hat B^2(R)\big)\,(\Delta p^2/p_H^2)\,(p_H^2R^2-I^2)\ \ge1$, with
$\hat B^2(R)=(\Delta p^2/p_H^2)\big(p_H^2R^2-I^2+R\big)$ and gross utility
$W(R)=R+F(\hat B^2(R))(p_H^2R^2-I^2)$.

\step{4}{(ii) Date-1 incentive constraint} By \eqref{t4e11-W} and Step 1, $W(0)=0+F(\hat B^2(0))\,\mathit{NPV}^2=0$, so
the entrepreneur works at date 1 iff
\begin{align}
  p_H^1W(R_b^1)+(1-p_H^1)W(0) &\ge p_L^1W(R_b^1)+(1-p_L^1)W(0)+B^1 \nonumber\\
  (p_H^1-p_L^1)W(R_b^1)-(p_H^1-p_L^1)W(0) &\ge B^1 \nonumber\\
  (\Delta p^1)\,W(R_b^1) &\ge B^1 \qquad (W(0)=0) \nonumber\\
  (\Delta p^1)\big[R_b^1+F\big(\hat B^2(R_b^1)\big)\mathit{NPV}^2\big] &\ge B^1, \tag{IC$^1$}\label{t4e11-ic}
\end{align}
the last line substituting \eqref{t4e11-W}.
Let $\tilde R_b^1$ be the smallest $R_b^1$ satisfying \eqref{t4e11-ic}. Whenever
$F(\hat B^2(\tilde R_b^1))>0$,
\begin{align}
  \tilde R_b^1 = \frac{B^1}{\Delta p^1}-F\big(\hat B^2(\tilde R_b^1)\big)\mathit{NPV}^2 < \frac{B^1}{\Delta p^1}. \label{t4e11-Rt}
\end{align}

\step{5}{(ii) Date-1 financing} The entrepreneur maximises $p_H^1W(R_b^1)-A^1$ subject to
\eqref{t4e11-ic} and $p_H^1(R^1-R_b^1)\ge I^1-A^1$. Since $W$ is increasing \eqref{t4e11-mu}, the
breakeven constraint binds, and financing is possible iff
\begin{align}
  p_H^1\big(R^1-\tilde R_b^1\big)\ \ge\ I^1-A^1. \label{t4e11-fin}
\end{align}
By \eqref{t4e11-Rt} the left-hand side exceeds the stand-alone pledgeable income
$p_H^1(R^1-B^1/\Delta p^1)$. So for $A^1$ with
$p_H^1(R^1-B^1/\Delta p^1)<I^1-A^1\le p_H^1(R^1-\tilde R_b^1)$, project 1 is funded only
because project 2 exists, although none of project 2's income is pledged at date 1.

\step{6}{(ii) An example} Project 2: $I^2=20$, $R^2=40$, $p_H^2=0.75$, $p_L^2=0.25$, $B^2\sim U[7,9]$,
so $\mathit{NPV}^2=10$ and $\Delta p^2\mathit{NPV}^2/p_H^2=20/3<7$ (assumption holds; also $p_L^2R^2+B^2<I^2$).
By \eqref{t4e11-Bhat}, $\hat B^2(R)=\tfrac23(10+R)$, which lies in $[7,9]$ for $R\in[0.5,3.5]$; there
$F=(\hat B^2-7)/2$ and \eqref{t4e11-W}--\eqref{t4e11-mu} give
\begin{align}
  W(R) &= R+\frac{\tfrac23(10+R)-7}{2}\,(10) \nonumber\\
       &= R+\frac{(20+2R-21)/3}{2}\,(10) = R+\frac{5(2R-1)}{3} \nonumber\\
       &= \frac{3R+10R-5}{3}=\frac{13R-5}{3}, \nonumber\\
  \mu &= 1+\Big(\frac12\Big)\Big(\frac{0.5}{0.75}\Big)(10)=1+\frac{10}{3}=\frac{13}{3}. \nonumber
\end{align}
Project 1: $I^1=4.5$, $R^1=10$, $p_H^1=0.5$, $p_L^1=0.25$, $B^1=0.5$, so $B^1/\Delta p^1=2$. Alone, the
pledgeable income is $(0.5)(10-2)=4$ and $\bar A^1=4.5-4=0.5$. With project 2, \eqref{t4e11-ic} at
equality reads $W(\tilde R_b^1)=B^1/\Delta p^1$, i.e.\ $(13\tilde R_b^1-5)/3=2$, so
$13\tilde R_b^1=6+5=11$, $\tilde R_b^1=11/13\approx0.846154$ (inside $[0.5,3.5]$, where the formula
for $W$ holds), and
\begin{align}
  p_H^1(R^1-\tilde R_b^1) = (0.5)\Big(10-\frac{11}{13}\Big)=\frac{119}{26}\approx 4.576923 > 4.5. \nonumber
\end{align}
The threshold falls from $0.5$ to $4.5-119/26=-1/13\approx-0.076923$: any $A^1\in[0,0.5)$ is now funded
(at $A^1=0$, $R_b^1=10-4.5/0.5=1\ge\tilde R_b^1$).

\answer{(ii)} Yes. Because project 2's rent is lost on failure, the date-1 reward needed for effort
falls to $\tilde R_b^1<B^1/\Delta p^1$ and date-1 pledgeable income rises to
$p_H^1(R^1-\tilde R_b^1)>p_H^1(R^1-B^1/\Delta p^1)$. Any $A^1$ with
$p_H^1(R^1-B^1/\Delta p^1)<I^1-A^1\le p_H^1(R^1-\tilde R_b^1)$ is financed only thanks to project 2
(in the example, $A^1\in[0,0.5)$).

\intuition Retained earnings are worth more than their face value because they are the net worth that
gets the next project financed. That future rent is lost if project 1 fails, so it acts as a reward for
date-1 effort: part of the incentive is paid with continuation value instead of cash, and the cash
saved becomes pledgeable to date-1 lenders.

\pitfalls
\begin{itemize}
  \item Valuing retained earnings at face ($\mu=1$). Each unit also raises the chance that project 2 is
        financed, worth $f\,(\Delta p^2/p_H^2)\,\mathit{NPV}^2$ more.
  \item Writing the date-1 incentive constraint with $R_b^1$ alone, $(\Delta p^1)R_b^1\ge B^1$: the
        reward for success is $W(R_b^1)$, which includes the continuation rent.
  \item Treating project 2's income as pledgeable at date 1. The effect works only through the
        entrepreneur's incentives; date-1 lenders are repaid out of $R^1$.
\end{itemize}

\bookans Matches the book's answer (p.~575): $\mu=1+f(\hat B^2)(d\hat B^2/dR_b^1)\mathit{NPV}^2$ and the
date-1 constraint $(\Delta p^1)[R_b^1+F(\hat B^2(R_b^1))\mathit{NPV}^2]\ge B^1$. The book writes the
constraint without $W(0)$; this is right because the assumption on $\underline B^2$ makes $W(0)=0$
(Step 1). The numerical example is this solution's.

\refbib Tirole \S3.2.2, pp.~116--117; \S3.4.2, p.~128 (shadow value of equity); answer p.~575.
%% END-FROM

% =============================================================================
\exsec{12. Continuous investment and decreasing returns to scale --- Tirole Exercise 3.5, p.~146}

\asked (i) How the investment the entrepreneur can carry out, $I(A)$, changes with her assets
$A$ when returns are decreasing, $R(I)$ concave. (ii) How the shadow value of assets
$v=dU_b^g/dA$ (the marginal value of one more unit of cash, in gross utility) changes with $A$.

\approach Variable-investment Holmstr\"om--Tirole model (Tirole \S3.4) with a concave return.
The incentive constraint still requires $R_b\ge BI/\Delta p$, so pledgeable income is
$p_H[R(I)-BI/\Delta p]$; competitive investors break even, so the entrepreneur maximizes the NPV
$p_HR(I)-I$ subject to the investors' breakeven constraint. Kuhn--Tucker tells which constraint
binds; implicit differentiation of the binding constraint gives $I'(A)$; the envelope theorem
gives $v$ as $1$ plus the multiplier of the breakeven constraint.

\notation
\begin{itemize}
  \item $R(I)$: income on success, $R'>0$, $R''<0$, $R'(0)>1/p_H>R'(\infty)$; $0$ on failure.
  \item $p_H$, $p_L$, $\Delta p$; $BI$: private benefit from misbehaving; $A$: cash; $R_b$: entrepreneur's reward on success.
  \item $I^*$: first-best scale, $p_HR'(I^*)=1$; $A^*$: cash needed to finance $I^*$.
  \item $U_b=p_HR(I)-I$: net utility (NPV); $U_b^g=A+U_b$: gross utility; $v=dU_b^g/dA$; $\mu$: multiplier of the breakeven constraint.
  \item $D(I)\equiv 1-p_HR'(I)+p_HB/\Delta p$.
\end{itemize}

\step{1}{Pledgeable income} The incentive constraint is $(\Delta p)R_b\ge BI$, so the smallest
stake is $R_b=BI/\Delta p$ and the most that can be promised to investors in expectation is
\begin{align}
  \mathcal P(I) &= p_H\Big[R(I)-\frac{BI}{\Delta p}\Big]. \label{r4e12-P}
\end{align}

\step{2}{The program} Investors break even, so the entrepreneur's net utility is the NPV and she
solves
\begin{align}
  \max_I\ & p_HR(I)-I \nonumber\\
  \text{s.t. }\ & p_H\Big[R(I)-\frac{BI}{\Delta p}\Big]\ge I-A. \tag{BB}
\end{align}

\step{3}{The unconstrained optimum} The objective is concave; its derivative $p_HR'(I)-1$ is
positive at $I=0$ (since $R'(0)>1/p_H$), negative for large $I$ (since $R'(\infty)<1/p_H$), and
decreasing (since $R''<0$). So it has a unique maximum $I^*$, given by $p_HR'(I^*)=1$. It is
feasible if (BB) holds at $I^*$, i.e.\ if
\begin{align}
  A\ge A^* &\equiv I^*-p_H\Big[R(I^*)-\frac{BI^*}{\Delta p}\Big]. \label{r4e12-Astar}
\end{align}
Then $I(A)=I^*$ and the multiplier is $\mu=0$.

\step{4}{The constrained region $A<A^*$} Define $g(I)=p_HR(I)-(p_HB/\Delta p+1)I+A$, so (BB) is
$g(I)\ge 0$. $g$ is concave and $g(0)=p_HR(0)+A\ge 0$ (with $R(0)\ge 0$), so the feasible set is an interval
$[0,\bar I(A)]$, where $\bar I(A)$ is the larger root of $g=0$, and $g'(\bar I)\le 0$ there, i.e.\
$D(\bar I)=-g'(\bar I)\ge 0$ (strictly, generically). When $A<A^*$, $I^*>\bar I(A)$; the concave
objective increases up to $I^*$, so on $[0,\bar I(A)]$ it is maximized at the boundary:
\begin{align}
  I(A)=\bar I(A):\qquad p_H\Big[R(I)-\frac{BI}{\Delta p}\Big] &= I-A. \label{r4e12-bind}
\end{align}
The Lagrangian $\mathcal L=p_HR(I)-I+\mu\big[p_HR(I)-p_HBI/\Delta p-I+A\big]$ has first-order condition
\begin{align}
  p_HR'(I)-1+\mu\Big[p_HR'(I)-\frac{p_HB}{\Delta p}-1\Big] &= 0 \nonumber\\
  \mu &= \frac{p_HR'(I)-1}{1+p_HB/\Delta p-p_HR'(I)} = \frac{p_HR'(I)-1}{D(I)} > 0, \label{r4e12-mu}
\end{align}
positive because $I<I^*$ makes $p_HR'(I)>1$.

\step{5}{(i) Slope and curvature of $I(A)$} Differentiate \eqref{r4e12-bind} with respect to $A$,
with $I=I(A)$:
\begin{align}
  p_HR'(I)\,I'-\frac{p_HB}{\Delta p}\,I' &= I'-1 \nonumber\\
  I'\Big[1-p_HR'(I)+\frac{p_HB}{\Delta p}\Big] &= 1 \nonumber\\
  I'(A) &= \frac{1}{D(I)} > 0. \label{r4e12-Ip}
\end{align}
Differentiate again, using $dD/dA=-p_HR''(I)\,I'$:
\begin{align}
  I''(A) &= -\frac{1}{D^2}\,\frac{dD}{dA} = \frac{p_HR''(I)\,I'(A)}{D(I)^2} < 0. \nonumber
\end{align}

\answer{(i)} For $A<A^*$, $I(A)$ is given by the binding breakeven constraint
\eqref{r4e12-bind}; it is increasing and concave, with slope $1/[1-p_HR'(I)+p_HB/\Delta p]$. For
$A\ge A^*$, $I(A)=I^*$ (flat).

\step{6}{(ii) The shadow value} For $A<A^*$, $U_b^g=A+p_HR(I(A))-I(A)$; differentiating and
substituting \eqref{r4e12-Ip},
\begin{align}
  v=\frac{dU_b^g}{dA} &= 1+\big[p_HR'(I)-1\big]I'(A) \nonumber\\
    &= 1+\frac{p_HR'(I)-1}{1-p_HR'(I)+p_HB/\Delta p} \nonumber\\
    &= \frac{1-p_HR'(I)+p_HB/\Delta p+p_HR'(I)-1}{1-p_HR'(I)+p_HB/\Delta p} \nonumber\\
    &= \frac{p_HB/\Delta p}{1-p_HR'(I(A))+p_HB/\Delta p}. \label{r4e12-v}
\end{align}
By \eqref{r4e12-mu}, $v=1+\mu$, which is the envelope theorem ($\partial\mathcal L/\partial A=\mu$
plus the direct $1$). For $A\ge A^*$, $v=1$ (net shadow value $0$). At $A=A^*$,
$p_HR'(I^*)=1$ and \eqref{r4e12-v} gives $(p_HB/\Delta p)/(p_HB/\Delta p)=1$: $v$ is continuous.

\step{7}{(ii) $v$ decreases with $A$} Differentiate \eqref{r4e12-v}:
\begin{align}
  \frac{dv}{dA} &= -\frac{p_HB/\Delta p}{D(I)^2}\,\frac{dD}{dA}
   = -\frac{p_HB/\Delta p}{D(I)^2}\big(-p_HR''(I)\big)I'(A) < 0, \nonumber
\end{align}
since $-p_HR''>0$ and $I'>0$.

\answer{(ii)} $v(A)=\dfrac{p_HB/\Delta p}{1-p_HR'(I(A))+p_HB/\Delta p}>1$ and strictly decreasing
for $A<A^*$; $v=1$ for $A\ge A^*$. (With constant returns, $R'=R$ and $v$ is the constant
(3.15).)

\emph{Numerical illustration.} $R(I)=4\sqrt I$, $p_H=0.5$, $B/\Delta p=3$ (e.g.\ $B=0.75$,
$\Delta p=0.25$). Then $R'(I)=2/\sqrt I$, so $R'(0)=\infty>2$ and $R'(\infty)=0<2$; $p_HR'(I)=1/\sqrt I$,
$I^*=1$; $\mathcal P(I)=0.5(4\sqrt I-3I)=2\sqrt I-1.5I$; $A^*=1-(2-1.5)=0.5$; $p_HB/\Delta p=1.5$.
For $A<0.5$, \eqref{r4e12-bind} reads $2\sqrt I-1.5I=I-A$, i.e.\ with $s=\sqrt I$,
$2.5s^2-2s-A=0$, larger root $s=[2+\sqrt{4+10A}]/5$:
\begin{center}
\begin{tabular}{cccccc}
\toprule
$A$ & $s=\sqrt{I}$ & $I(A)$ & $D=1-1/s+1.5$ & $I'=1/D$ & $v=1.5/D$\\
\midrule
$0$   & $0.8$      & $0.64$     & $1.25$     & $0.8$      & $1.2$\\
$0.2$ & $0.889898$ & $0.791918$ & $1.376276$ & $0.726599$ & $1.089898$\\
$0.5$ & $1$        & $1$        & $1.5$      & $0.666667$ & $1$\\
\bottomrule
\end{tabular}
\end{center}
(at $A=0.2$, $1/s=p_HR'=1.123724$). Investment rises at a falling rate and $v$ falls from $1.2$
to $1$ at $A^*=0.5$.

\intuition With decreasing returns, the first units of investment are very productive and
pledgeable income is large relative to cost, so even a poor entrepreneur invests something.
Each extra unit of cash relaxes the breakeven constraint and lets her expand, but the units she
adds are less and less productive, so a unit of cash is worth less the richer she is. Once she
can fund the first-best scale, cash is worth exactly its face value. With constant returns
(Tirole \S3.4) the shadow value is instead a constant above $1$.
\pitfalls
\begin{itemize}
  \item Confusing gross and net shadow values: the derivative of the net utility $U_b$ is
        $v-1$, which is $0$ (not $1$) once $I^*$ is affordable.
  \item Picking the smaller root of the breakeven condition: the entrepreneur wants the
        largest feasible scale below $I^*$.
  \item Writing $I=kA$ as in the constant-returns model; with concave $R$ the relation is
        nonlinear and the ``multiplier'' $I/A$ changes with $A$.
\end{itemize}
\bookans Matches the book's answer (p.~573): (BB) binds when $A$ is small, $\mu=0$ when $I^*$ is
feasible, and the shadow value decreases with $A$. The book's displayed
$dU_b/dA=1/\{(p_HB/\Delta p)/(p_HR'-1)-1\}$ is the derivative of \emph{net} utility; it equals
$(p_HR'-1)/(1-p_HR'+p_HB/\Delta p)=v-1$ in our notation, which is the same result.
\refbib{Tirole Ch.~3, \S3.4.1--3.4.2 ``The Continuous-Investment Model'', ``The Lenders' Credit Analysis'', eq.~(3.15), pp.~127--128.}

% =============================================================================
\exsec{13. Diversification with variable-investment projects --- Tirole Exercise 4.10, pp.~191--192}

\asked (i) If each of two independent variable-investment projects is financed on its own
(project finance), what utility does the entrepreneur get, and does having two projects help?
(ii) What utility does she get if the incomes of the two projects can be pledged jointly
(cross-pledging)?

\approach Project finance is just the model of Tirole \S3.4 twice, so the equity multiplier
$1/(1-\rho_0)$ applies to each project and total utility depends only on total cash. Under
cross-pledging (Tirole \S4.2.1.2) the entrepreneur's reward can depend on both outcomes; the
cheapest way to provide incentives is to pay only when both projects succeed (the outcome most
indicative of effort), which lowers the expected rent per unit of investment and raises
pledgeable income from $\rho_0$ to $\rho_0'$.

\notation
\begin{itemize}
  \item $I^i$: investment in project $i\in\{1,2\}$; $I=I^1+I^2$; $A$: total cash; $A_i$: cash put into project $i$.
  \item Project $i$ yields $RI^i$ on success, $0$ on failure; $p_H$, $p_L=p_H-\Delta p$; private benefit $BI^i$ from misbehaving on $i$.
  \item $\rho_1=p_HR>1>\rho_0=p_H(R-B/\Delta p)$; $\rho_0'=p_H\big(R-\frac{p_H}{p_H+p_L}\frac{B}{\Delta p}\big)<1$.
  \item $R_b^i$: reward on success of project $i$ (project finance); $R_k$, $k\in\{0,1,2\}$: reward when $k$ projects succeed (cross-pledging); $m$: expected reward.
  \item $\pi_k^H$, $\pi_k^L$: probability of $k$ successes if she works on both / shirks on both.
\end{itemize}

\step{1}{(i) Project finance} Project $i$ is financed alone with cash $A_i$, $A_1+A_2=A$. Its
constraints are those of Exercise~7(i) with $(I^i,A_i)$ in place of $(I,A)$, so by \eqref{r4e7-k}
and \eqref{r4e7-Ub},
\begin{align}
  I^i &= \frac{A_i}{1-\rho_0}, \qquad U_b^i=(\rho_1-1)I^i. \nonumber
\end{align}
Adding the two projects,
\begin{align}
  U_b &= (\rho_1-1)\frac{A_1}{1-\rho_0}+(\rho_1-1)\frac{A_2}{1-\rho_0} \nonumber\\
      &= \frac{\rho_1-1}{1-\rho_0}(A_1+A_2) = \frac{\rho_1-1}{1-\rho_0}A. \label{r4e13-pf}
\end{align}

\answer{(i)} $U_b=\dfrac{\rho_1-1}{1-\rho_0}A$, whatever the split of $A$: the same as with a
single project, so there is no benefit from a second project under project finance.

\step{2}{(ii) Outcome probabilities} With cross-pledging the reward depends on the number of
successes. With independent projects,
\begin{align}
  \pi_2^H &= p_H^2, & \pi_1^H &= 2p_H(1-p_H), & \pi_0^H &= (1-p_H)^2, \nonumber\\
  \pi_2^L &= p_L^2, & \pi_1^L &= 2p_L(1-p_L), & \pi_0^L &= (1-p_L)^2, \nonumber
\end{align}
and if she shirks on project $i$ only, both succeed with probability $p_Hp_L$.

\step{3}{(ii) Reward only joint success} We find the cheapest scheme that deters shirking on both
projects (private benefit $B(I^1+I^2)=BI$). Define the likelihood ratios
$\ell_k=(\pi_k^H-\pi_k^L)/\pi_k^H$:
\begin{align}
  \ell_2 &= 1-\frac{p_L^2}{p_H^2}, \qquad
  \ell_1 = 1-\frac{p_L(1-p_L)}{p_H(1-p_H)}, \qquad
  \ell_0 = 1-\frac{(1-p_L)^2}{(1-p_H)^2}<0. \nonumber
\end{align}
$\ell_1<\ell_2$ because
\begin{align}
  \frac{p_L(1-p_L)}{p_H(1-p_H)}>\frac{p_L^2}{p_H^2}
  &\iff \frac{p_H(1-p_L)}{1-p_H}>p_L \qquad(\text{multiplying by }p_H^2/p_L>0) \nonumber\\
  &\iff p_H-p_Hp_L>p_L-p_Hp_L \iff p_H>p_L. \nonumber
\end{align}
The incentive constraint is $\sum_k(\pi_k^H-\pi_k^L)R_k\ge BI$. Since $\pi_k^H-\pi_k^L=\ell_k\pi_k^H\le\ell_2\pi_k^H$
and $R_k\ge 0$,
\begin{align}
  BI \le \sum_k\ell_k\pi_k^HR_k \le \ell_2\sum_k\pi_k^HR_k = \ell_2\,m
  \;\Rightarrow\; m\ge\frac{BI}{\ell_2}=\frac{p_H^2BI}{p_H^2-p_L^2}=\frac{p_H^2BI}{(p_H+p_L)\Delta p}, \label{r4e13-m}
\end{align}
with equality when $R_1=R_0=0$ and
\begin{align}
  R_2 &= \frac{BI}{p_H^2-p_L^2}=\frac{BI}{(p_H+p_L)\Delta p}. \label{r4e13-R2}
\end{align}

\step{4}{(ii) The single-shirking constraint also holds} Shirking on project $i$ only gives
$p_Hp_LR_2+BI^i$, so it is deterred if
\begin{align}
  p_H^2R_2 &\ge p_Hp_LR_2+BI^i \nonumber\\
  p_H(\Delta p)R_2 &\ge BI^i. \nonumber
\end{align}
Substituting \eqref{r4e13-R2}: $p_H(\Delta p)\frac{BI}{(p_H+p_L)\Delta p}=\frac{p_H}{p_H+p_L}BI\ge BI^i$,
i.e.\ $I^i\le\frac{p_H}{p_H+p_L}I$. Nothing else depends on the split, so take $I^1=I^2=I/2$:
the condition becomes $\frac12\le\frac{p_H}{p_H+p_L}$, i.e.\ $p_L\le p_H$, which holds. Hence
\eqref{r4e13-R2} is optimal.

\step{5}{(ii) Financing capacity} Investors break even (binding, as in Exercise~7) on expected
income $p_HRI$ minus the expected reward \eqref{r4e13-m}:
\begin{align}
  p_HRI-\frac{p_H^2BI}{(p_H+p_L)\Delta p} &\ge I-A \nonumber\\
  p_H\Big(R-\frac{p_H}{p_H+p_L}\frac{B}{\Delta p}\Big)I &\ge I-A \nonumber\\
  \rho_0'I &\ge I-A \nonumber\\
  I &\le \frac{A}{1-\rho_0'} \qquad(\text{dividing by }1-\rho_0'>0). \label{r4e13-I}
\end{align}

\step{6}{(ii) Utility} As in \eqref{r4e7-Ub}, the borrower receives the NPV, $U_b=(\rho_1-1)I$,
and invests the maximum \eqref{r4e13-I}:
\begin{align}
  U_b &= \frac{\rho_1-1}{1-\rho_0'}A. \label{r4e13-cp}
\end{align}
The gain over project finance comes from
\begin{align}
  \rho_0'-\rho_0 &= p_H\frac{B}{\Delta p}\Big(1-\frac{p_H}{p_H+p_L}\Big)
                  = p_H\frac{B}{\Delta p}\cdot\frac{p_L}{p_H+p_L}>0, \nonumber
\end{align}
where $p_L/(p_H+p_L)=d_2$ is Tirole's diversification measure.

\answer{(ii)} $U_b=\dfrac{\rho_1-1}{1-\rho_0'}A>\dfrac{\rho_1-1}{1-\rho_0}A$: cross-pledging
raises the borrowing capacity and hence the NPV.

\emph{Numerical illustration.} $p_H=0.8$, $p_L=0.4$, $R=1.5$, $B=0.3$, $A=10$: $\rho_1=1.2$,
$\rho_0=0.6$, $p_H/(p_H+p_L)=2/3$, $\rho_0'=0.8(1.5-\tfrac23\cdot0.75)=0.8(1.5-0.5)=0.8$.
Project finance: $I=10/0.4=25$ ($12.5$ per project with $A_i=5$; $R_b^i=0.3(12.5)/0.4=9.375$;
pledgeable $0.8(1.5\cdot12.5-9.375)=7.5=12.5-5$), $U_b=0.2(25)=5$. Cross-pledging:
$I=10/0.2=50$, $U_b=0.2(50)=10$; $R_2=0.3(50)/(1.2\cdot0.4)=31.25$, $m=0.64(31.25)=20$, investors
get $0.8(1.5)(50)-20=60-20=40=I-A$; both-shirk constraint $(0.64-0.16)(31.25)=15=BI$;
single-shirk $0.8(0.4)(31.25)=10\ge BI^i=0.3(25)=7.5$. Here $\ell_2=1-0.25=0.75$ and
$\ell_1=1-0.24/0.16=-0.5$.

\intuition Under project finance, a failure on one project costs the entrepreneur nothing on the
other, so each project carries its own full incentive rent. Pooling lets the income of a
successful project serve as collateral for the other: she is paid only if both succeed, an
outcome that is much less likely when she shirks ($p_L^2$ against $p_H^2$), so a smaller
expected reward deters shirking. The rent per unit of investment falls, pledgeable income per
unit rises from $\rho_0$ to $\rho_0'$, and with variable scale this translates into a larger
investment and a larger NPV, not just easier access as with fixed projects.
\pitfalls
\begin{itemize}
  \item Paying a reward when exactly one project succeeds ($R_1>0$): that is project finance in
        disguise and wastes incentive money.
  \item Checking only one of the two deviations (shirk on both, shirk on one); both must hold,
        and the equal split $I^1=I^2$ is what makes the second one slack.
  \item Expecting a higher NPV per unit of investment: diversification raises the \emph{scale},
        the NPV per unit stays $\rho_1-1$.
\end{itemize}
\bookans Matches the book's answer (p.~587): $U_b=(\rho_1-1)A/(1-\rho_0)$ under project finance,
with no benefit from a second project; under cross-pledging $R_2>0=R_1=R_0$, $I^1=I^2=I/2$, and
$U_b=(\rho_1-1)A/(1-\rho_0')$.
\refbib{Tirole Ch.~3, \S3.4.2, pp.~127--128; Ch.~4, \S4.2.1 ``The Benefits of Diversification: The Case of Two Projects'', pp.~158--160.}

% =============================================================================
\exsec{14. Variable effort, marginal value of net worth, pooling of equity --- Tirole Exercise 3.20, p.~152}

\asked (i) Check that without borrowing the entrepreneur's net utility is $\tfrac12R^2-I$. (ii)
Find the cash threshold $\bar A$ below which the project is not funded, and show that the
marginal value of net worth $V'(A)$ is positive, zero at $A=I$, and decreasing. (iii) Two
entrepreneurs, one of whom will turn out to need $I_H$ (and could not raise it alone) and the
other $I_L$: should they agree in advance to pool their cash so that the lucky one subsidizes
the unlucky one?

\approach Continuous effort replaces the binary choice: with stake $R_b$ the entrepreneur picks
$p=R_b$, so pledgeable income $R_b(R-R_b)$ is hump-shaped (a ``Laffer curve'' in the
entrepreneur's stake). Competitive investors break even, so her utility is the NPV, which falls as
her stake is diluted. Implicit differentiation of the breakeven condition gives $V'$ and $V''$.
Part (iii) applies concavity of the value of net worth: moving cash where it is worth most
(the firm that cannot otherwise invest) can pay, at the cost of weaker incentives in the other.

\notation
\begin{itemize}
  \item $I>A$: fixed investment; $A$: cash; $R\in(\sqrt{2I},1)$: income on success; $p$: success probability, effort cost $\tfrac12p^2$.
  \item $R_b$: entrepreneur's reward on success; $\mathcal P(R_b)=p(R-R_b)$: pledgeable income; $W(R_b)=pR-\tfrac12p^2-I$: NPV at the effort induced by $R_b$.
  \item $V(A)$: NPV (= borrower's net utility) when financed with cash $A$; $V^*=\tfrac12R^2-I$; $A_1=I-\tfrac14R^2$; $A_2$: cash at which $V=0$.
  \item (iii): $I_L<I_H$; $\bar I=(I_L+I_H)/2$; $V_I(\cdot)$: the function $V$ for investment $I$; $A_L$, $A_H$: cash allocated to the lucky / unlucky project.
\end{itemize}

\step{1}{(i) No borrowing} With $A\ge I$ the entrepreneur is residual claimant, $R_b=R$, and she
chooses $p$ to maximize $pR-\tfrac12p^2-I$. The first-order condition is $R-p=0$, so $p=R$
(interior because $R<1$), and
\begin{align}
  U_b &= (R)R-\tfrac12(R)^2-I = \tfrac12R^2-I = V^*, \nonumber
\end{align}
positive because $R>\sqrt{2I}\iff R^2>2I$.

\answer{(i)} $U_b=V^*=\tfrac12R^2-I>0$, independent of $A$ for $A\ge I$.

\step{2}{(ii) Effort and pledgeable income} With stake $R_b$, she maximizes $pR_b-\tfrac12p^2$:
the first-order condition $R_b-p=0$ gives $p=R_b$. Investors receive $R-R_b$ with probability $p$:
\begin{align}
  \mathcal P(R_b) &= R_b(R-R_b), \qquad \mathcal P'(R_b)=R-2R_b, \label{r4e14-P}
\end{align}
so $\mathcal P$ is maximized at $R_b=\tfrac12R$, where $\mathcal P=\tfrac12R\cdot\tfrac12R=\tfrac14R^2$.

\step{3}{(ii) Only $R_b\ge\tfrac12R$ matters} $\mathcal P$ is symmetric around $\tfrac12R$:
$\mathcal P(R-R_b)=(R-R_b)R_b=\mathcal P(R_b)$. The NPV at the induced effort is
$W(R_b)=R_bR-\tfrac12R_b^2-I$, with $W'(R_b)=R-R_b>0$ for $R_b<R$. So for any $R_b<\tfrac12R$, the
stake $R-R_b>R_b$ yields the same pledgeable income and a higher NPV: we can restrict attention
to $R_b\in[\tfrac12R,R]$.

\step{4}{(ii) The borrower's utility is the NPV} Her net utility is $pR_b-\tfrac12p^2-A=\tfrac12R_b^2-A$.
With binding breakeven $R_b(R-R_b)=I-A$, i.e.\ $A=I-R_bR+R_b^2$,
\begin{align}
  U_b &= \tfrac12R_b^2-(I-R_bR+R_b^2) = R_bR-\tfrac12R_b^2-I = W(R_b). \nonumber
\end{align}
She wants the largest $R_b$ compatible with breakeven, the larger root of
$R_b^2-RR_b+(I-A)=0$:
\begin{align}
  R_b(A) &= \frac{R+\sqrt{R^2-4(I-A)}}{2}, \qquad V(A)=W\big(R_b(A)\big). \label{r4e14-Rb}
\end{align}

\step{5}{(ii) The threshold} A real root requires $R^2-4(I-A)\ge 0$, i.e.\
$A\ge A_1=I-\tfrac14R^2$ (the maximum pledgeable income covers $I-A$). The NPV must also be
nonnegative. At $A_1$, $R_b=\tfrac12R$ and
\begin{align}
  V(A_1) &= \tfrac12R\cdot R-\tfrac12\big(\tfrac12R\big)^2-I = \tfrac12R^2-\tfrac18R^2-I = \tfrac38R^2-I. \nonumber
\end{align}
If $\tfrac38R^2\ge I$, the threshold is $A_1$. Otherwise $V(A_1)<0$ and the threshold is $A_2>A_1$
with $V(A_2)=0$: $R_bR-\tfrac12R_b^2=I\iff R_b^2-2RR_b+2I=0$, root $R_b=R-\sqrt{R^2-2I}$ (which is
$\ge\tfrac12R$ exactly when $I\ge\tfrac38R^2$); writing $s=\sqrt{R^2-2I}$,
\begin{align}
  A_2 &= I-R_b(R-R_b) = I-(R-s)s = I-Rs+(R^2-2I) = R^2-I-R\sqrt{R^2-2I}. \nonumber
\end{align}

\answer{(ii-a)} $\bar A=\max\{A_1,A_2\}$ with $A_1=I-\tfrac14R^2$ and
$A_2=R^2-I-R\sqrt{R^2-2I}$ (relevant only if $I>\tfrac38R^2$).

\step{6}{(ii) $V'(A)$} By the chain rule, $V'(A)=W'(R_b)\,R_b'(A)=(R-R_b)R_b'(A)$. Differentiate
the breakeven identity $R_b(R-R_b)=I-A$:
\begin{align}
  (R-2R_b)R_b' &= -1 \nonumber\\
  R_b'(A) &= \frac{1}{2R_b-R} > 0 \qquad(R_b>\tfrac12R), \nonumber\\
  V'(A) &= \frac{R-R_b}{2R_b-R} > 0 \qquad(\tfrac12R<R_b<R). \label{r4e14-Vp}
\end{align}
At $A=I$, \eqref{r4e14-Rb} gives $R_b=\frac{R+\sqrt{R^2}}{2}=R$, so $V'(I)=0$; for $A>I$, $V=V^*$ is
flat, so $V'$ is continuous at $I$.

\step{7}{(ii) $V''(A)$} Let $h(R_b)=(R-R_b)/(2R_b-R)$. By the quotient rule,
\begin{align}
  h'(R_b) &= \frac{-(2R_b-R)-2(R-R_b)}{(2R_b-R)^2} = \frac{-2R_b+R-2R+2R_b}{(2R_b-R)^2} = \frac{-R}{(2R_b-R)^2}, \nonumber\\
  V''(A) &= h'(R_b)\,R_b'(A) = -\frac{R}{(2R_b-R)^3} < 0. \label{r4e14-Vpp}
\end{align}

\answer{(ii-b)} $V'(A)=\dfrac{R-R_b(A)}{2R_b(A)-R}>0$ for $\bar A\le A<I$, $V'(I)=0$, and
$V''(A)=-R/(2R_b-R)^3<0$.

\emph{Numerical illustration.} $R=0.9$, $I=0.3$ ($\sqrt{2I}=0.774597<0.9<1$): $V^*=0.405-0.3=0.105$;
$\tfrac14R^2=0.2025$, $A_1=0.3-0.2025=0.0975$; $\tfrac38R^2=0.30375\ge 0.3$, so $\bar A=0.0975$ and
$V(\bar A)=0.30375-0.3=0.00375$. At $A=0.2$: $\sqrt{0.81-0.4}=0.640312$,
$R_b=(0.9+0.640312)/2=0.7701562$, $V=0.6931406-0.2965703-0.3=0.0965703$ (also
$\tfrac12R_b^2-A=0.2965703-0.2=0.0965703$), $V'=(0.9-0.7701562)/0.640312=0.129844/0.640312=0.202782$,
$V''=-0.9/0.640312^3=-3.4282$.

\step{8}{(iii) Without pooling} The lucky entrepreneur ($I_L$) is funded: $I_L-A<\tfrac14R^2$, and
$\tfrac38R^2>I_H>I_L$ makes the NPV positive. The unlucky one ($I_H$) is not: even the maximum
pledgeable income $\tfrac14R^2$ is below $I_H-A$. Each is lucky with probability $\tfrac12$, so
ex ante each expects $\tfrac12V_{I_L}(A)$.

\step{9}{(iii) Pooling with both projects funded: equal stakes} With cash $2A$ pooled, rewards
$R_b^L$, $R_b^H$, and investors breaking even overall, total NPV is
$\sum_i\big[R_b^iR-\tfrac12(R_b^i)^2\big]-I_L-I_H$. The Lagrangian is
\begin{align}
  \mathcal L &= \sum_{i\in\{L,H\}}\Big[R_b^iR-\tfrac12(R_b^i)^2\Big]
     +\lambda\Big[\sum_iR_b^i(R-R_b^i)-(I_L+I_H-2A)\Big], \nonumber
\end{align}
and the first-order condition for each $i$,
\begin{align}
  R-R_b^i+\lambda(R-2R_b^i) &= 0 \;\Rightarrow\; R_b^i=\frac{R(1+\lambda)}{1+2\lambda}, \nonumber
\end{align}
is the same for both: $R_b^L=R_b^H=R_b$. The breakeven constraint becomes
$2R_b(R-R_b)=I_L+I_H-2A$, i.e.\ $R_b(R-R_b)=\bar I-A$, feasible because
$I_L+I_H-2A<\tfrac12R^2\iff\bar I-A<\tfrac14R^2$. Each entrepreneur's (equal) net utility is
\begin{align}
  V_{\bar I}(A) &= R_bR-\tfrac12R_b^2-\bar I, \label{r4e14-both}
\end{align}
positive because $R_b\ge\tfrac12R$ gives $R_bR-\tfrac12R_b^2\ge\tfrac38R^2>I_H>\bar I$. Each project's
own breakeven $R_b(R-R_b)=I_i-A_i$ fixes the split: $A_H-A_L=I_H-I_L$ and $A_L+A_H=2A$, so
\begin{align}
  A_L &= A-\tfrac12(I_H-I_L), \qquad A_H=A+\tfrac12(I_H-I_L): \nonumber
\end{align}
the lucky firm transfers $\tfrac12(I_H-I_L)$ of net worth to the unlucky one. (Note
$V_{\bar I}(A)=V_{I_L}(A_L)-\tfrac12(I_H-I_L)$: the lucky firm's NPV minus half the extra cost.)

\step{10}{(iii) Pooling with only the lucky project funded} The lucky firm can use up to $2A$
(the rest is returned), so total NPV is $V_{I_L}(\min\{2A,I_L\})$ and each expects
$\tfrac12V_{I_L}(\min\{2A,I_L\})>\tfrac12V_{I_L}(A)$, because $V_{I_L}$ is strictly increasing
below $I_L$ (Step~6) and $A<I_L$. Pooling always beats no pooling.

\step{11}{(iii) Cross-subsidize or not} The entrepreneurs commit to fund both projects if and
only if
\begin{align}
  V_{\bar I}(A) &> \tfrac12V_{I_L}\big(\min\{2A,I_L\}\big). \label{r4e14-cond}
\end{align}
Raising $I_H$ lowers the left side (higher average cost, more dilution of $R_b$) and leaves the
right side unchanged, so \eqref{r4e14-cond} holds for $I_H$ below a threshold. Illustration:
$R=0.9$, $A=0.08$, $I_L=0.25$ ($I_L-A=0.17<0.2025$):
\begin{itemize}
  \item no pooling: $R_b=(0.9+\sqrt{0.81-0.68})/2=(0.9+0.360555)/2=0.6302776$,
        $V_{I_L}(0.08)=0.5672498-0.1986249-0.25=0.1186249$, ex ante $0.059312$;
  \item pooling, lucky only ($2A=0.16$, $I_L-2A=0.09$): $R_b=0.7854102$,
        $V_{I_L}(0.16)=0.7068692-0.3084346-0.25=0.1484346$, ex ante $0.074217$;
  \item pooling, both, $I_H=0.285$: $\bar I=0.2675$, $\bar I-A=0.1875$, $R_b=0.5724745$,
        $V_{\bar I}=0.5152270-0.1638635-0.2675=0.0838635>0.074217$: cross-subsidize
        (transfer $\tfrac12(0.285-0.25)=0.0175$);
  \item pooling, both, $I_H=0.30$: $\bar I=0.275$, $\bar I-A=0.195$, $R_b=0.5366025$,
        $V_{\bar I}=0.4829423-0.1439711-0.275=0.0639711<0.074217$: pool, but fund only the lucky firm.
\end{itemize}
Both values of $I_H$ satisfy the exercise's assumptions ($I_H-A>0.2025$, $I_H<0.30375$,
$I_L+I_H-2A<0.405$).

\answer{(iii)} Yes, they always want to pool. They commit to the lucky firm cross-subsidizing the
unlucky one (equal stakes, transfer $\tfrac12(I_H-I_L)$) when \eqref{r4e14-cond} holds, which is
the case for $I_H$ below a threshold; for higher $I_H$ the pool funds only the lucky project.

\intuition With continuous effort, every unit of outside money dilutes the entrepreneur's stake
and lowers effort, so the value of net worth is smooth: very high near the funding threshold,
zero once she is fully self-financed, and decreasing in between. Concavity is the reason to
share: a unit of net worth is worth more where it is scarce. Pooling at least lets the lucky
firm use both entrepreneurs' cash; funding both projects as well brings a second NPV but forces
both stakes down to the common $R_b$, and when the unlucky project is too expensive that
dilution costs more than the extra project brings.
\pitfalls
\begin{itemize}
  \item Dropping the effort cost from the NPV ($R_bR-I$ instead of $R_bR-\tfrac12R_b^2-I$): the
        borrower's utility is net of the private cost of effort.
  \item Taking the smaller root of $R_b(R-R_b)=I-A$; it gives the same pledgeable income with a
        lower NPV.
  \item In (iii), comparing the lucky firm's NPV $V_{I_L}(A_L)$ with the alternatives instead of
        each entrepreneur's expected value $V_{\bar I}(A)$, which is lower by $\tfrac12(I_H-I_L)$.
\end{itemize}
\bookans (i) and (ii) match the book's answer (p.~580): $p=R_b$, $\mathcal P=R_b(R-R_b)$,
$A_1=I-\tfrac14R^2$, $\bar A=\max\{A_1,A_2\}$, $V'(A)=[R-R_b(A)]\,dR_b/dA>0$, $V'(I)=0$, $V''<0$.
(iii) matches the book's conclusions (pp.~580--581): equal stakes, pooling always optimal,
cross-subsidization from lucky to unlucky for $I_H$ below a threshold. Two differences in the
book's displays: its $V(A-\tfrac12(I_H-I_L))$ is the lucky firm's NPV, while the ex ante value per
entrepreneur subtracts $\tfrac12(I_H-I_L)$ (\eqref{r4e14-both}); and it writes
$\tfrac12V(\max(2A,I_L))$ where $\min$ is meant (with $\max$ the value would always be $V^*$).
\refbib{Tirole Ch.~3, \S3.2.1--3.2.2, pp.~115--117 (footnote~3 points to this exercise for continuous effort). Cestone and Fumagalli (2005), cited in the statement.}

% =============================================================================
\exsec{15. Hedging or gambling on net worth? --- Tirole Exercise 3.21, pp.~152--153}

\asked Before knowing her date-1 cash $A_0+\varepsilon$, the entrepreneur can hedge for free
(locking in $A_0$) or not (``gamble''). Show which she prefers in five settings: (i) fixed
investment, binary effort: hedge if $A_0\ge\bar A$, gamble if $A_0<\bar A$; (ii) fixed
investment, continuous effort: hedge; (iii) variable investment with private benefit $B(I)$:
hedge if $B$ is convex, indifferent if linear, gamble if concave; (iv) variable investment and
unobservable income: hedge; (v) with a date-0 contract with investors, insulating investment from
$\varepsilon$ (full hedging) is optimal even where gambling was.

\approach Froot--Scharfstein--Stein logic (Tirole \S5.4): a risk-neutral entrepreneur hedges if
and only if the date-1 value of cash, $U(A)$, is concave, because hedging gives $U(A_0)$ and
gambling $\mathbb{E}[U(A_0+\varepsilon)]$, with $\mathbb{E}[A_0+\varepsilon]=A_0$ (Jensen). Each
part derives $U$ from a model seen before: the step at $\bar A$ (Tirole \S3.2), the concave $V$
of Exercise~14, the investment rule $I(A)$ of the variable model with nonlinear $B(I)$, and the
self-financed technology $R(A)$. Part (v) is liquidity management: contracting at date~0 removes
the reason to gamble.

\notation
\begin{itemize}
  \item $A_0$: expected date-1 cash; $\varepsilon\in[\underline\varepsilon,\bar\varepsilon]$, $\mathbb{E}[\varepsilon]=0$, $A_0+\underline\varepsilon\ge 0$; $F(x)=\Pr(\varepsilon<x)$.
  \item $U(A)$: date-1 gross utility with cash $A$; $U^h=U(A_0)$ (hedge), $U^g=\mathbb{E}[U(A_0+\varepsilon)]$ (gamble).
  \item (i): $\bar A$ from $p_H(R-B/\Delta p)=I-\bar A$; $V=p_HR-I$. (ii): as Exercise~14. (iii): $\rho_1=p_HR$; $B(I)$, $B'>0$; $I(A)$. (iv): $R(I)$ concave, unobservable.
\end{itemize}

\step{0}{The comparison} Because $\mathbb{E}[A_0+\varepsilon]=A_0$, she hedges if
$U(A_0)>\mathbb{E}[U(A_0+\varepsilon)]$, which holds for every nondegenerate $\varepsilon$ when $U$
is strictly concave, and the reverse when $U$ is strictly convex (Jensen's inequality).

\step{1}{(i) Date-1 utility} From (3.6), a borrower with cash $A\ge\bar A$ is funded and gets the
NPV $V$; below $\bar A$ she is not funded and keeps her cash:
\begin{align}
  U(A) &= A+V\cdot\mathbf 1\{A\ge\bar A\}. \nonumber
\end{align}

\step{2}{(i) $A_0\ge\bar A$} Hedging gives $U^h=A_0+V$. Gambling gives
\begin{align}
  U^g &= \mathbb{E}[A_0+\varepsilon]+V\Pr(A_0+\varepsilon\ge\bar A) \nonumber\\
      &= A_0+V\big[1-F(\bar A-A_0)\big] < A_0+V = U^h, \nonumber
\end{align}
because a wide support ($\underline\varepsilon<\bar A-A_0$) makes $F(\bar A-A_0)>0$. She hedges.

\step{3}{(i) $A_0<\bar A$} Hedging gives $U^h=A_0$ (no funding). Gambling gives
$U^g=A_0+V[1-F(\bar A-A_0)]>A_0$, because a wide support ($\bar\varepsilon>\bar A-A_0$) gives a
positive chance of reaching $\bar A$. She gambles.

\emph{Numerical illustration.} $I=100$, $R=150$, $p_H=0.8$, $p_L=0.4$, $B=30$ ($\bar A=40$,
$V=20$, from Exercise~8), $\varepsilon$ uniform on $[-20,20]$. $A_0=50$: $U^h=70$;
$\Pr(\varepsilon\ge-10)=0.75$, $U^g=50+20(0.75)=65<70$: hedge. $A_0=30$: $U^h=30$;
$\Pr(\varepsilon\ge10)=0.25$, $U^g=30+20(0.25)=35>30$: gamble ($A_0+\underline\varepsilon=10\ge0$).

\answer{(i)} Hedge if $A_0\ge\bar A$ (gambling only risks losing funding); gamble if
$A_0<\bar A$ (gambling gives a chance of funding).

\step{4}{(ii) Breakeven and NPV in terms of $R_b$} As in Exercise~14, $p=R_b$; the investors'
breakeven condition is $R_b(R-R_b)\ge I-A$ and the NPV is
\begin{align}
  W(R_b) &= R_bR-\tfrac12R_b^2-I, \nonumber
\end{align}
with $R_b\in[\tfrac12R,R]$ without loss of generality (Exercise~14, Step~3). Since
$I-A_0<\tfrac14R^2$ and the support is small, she is funded in every state, and
$U(A)=A+V(A)$ with $V(A)=W(R_b(A))$.

\step{5}{(ii) Concavity} By \eqref{r4e14-Vp}--\eqref{r4e14-Vpp}, $V'>0$ and $V''<0$ for $A<I$;
for $A\ge I$, $V=V^*$ is flat and $V'(I)=0$, so $V$ is concave on the whole range. $U$ is then
strictly concave where $A<I$, and $U^h=A_0+V(A_0)>A_0+\mathbb{E}[V(A_0+\varepsilon)]=U^g$.

\emph{Numerical illustration.} $R=0.9$, $I=0.3$, $A_0=0.2$, $\varepsilon=\pm0.05$ with
probability $\tfrac12$ ($I-A=0.15<0.2025$ in the bad state). $V(0.15)$: $R_b=(0.9+0.458258)/2=0.6791288$,
$V=0.6112159-0.2306080-0.3=0.0806080$. $V(0.25)$: $R_b=(0.9+0.781025)/2=0.8405125$,
$V=0.7564612-0.3532306-0.3=0.1032306$. Then $\mathbb{E}[V]=(0.0806080+0.1032306)/2=0.0919193<V(0.2)=0.0965703$.

\answer{(ii)} She hedges: with continuous effort the value of net worth is concave, so risk
in net worth lowers expected utility.

\step{6}{(iii) The investment rule} With investment constrained by pledgeable income and good
behaviour induced, investment solves the binding breakeven condition
\begin{align}
  p_H\Big[RI-\frac{B(I)}{\Delta p}\Big] &= I-A
  \iff \phi(I)\equiv(1-\rho_1)I+\frac{p_H}{\Delta p}B(I)=A. \label{r4e15-phi}
\end{align}
``Constrained'' means the marginal pledgeable income is below $1$ at the solution:
$\phi'(I)=1-\rho_1+(p_H/\Delta p)B'(I)>0$. Differentiating $\phi(I(A))=A$ once and twice,
\begin{align}
  \phi'(I)\,I' &= 1 \;\Rightarrow\; I'(A)=\frac{1}{\phi'(I)}>0, \nonumber\\
  I''(A) &= -\frac{\phi''(I)\,I'}{\phi'(I)^2} = -\frac{(p_H/\Delta p)B''(I)}{\phi'(I)^3}. \nonumber
\end{align}
So $I$ is concave if $B''>0$, linear if $B''=0$, convex if $B''<0$.

\step{7}{(iii) Utility} Investors break even, so $U(A)=A+(\rho_1-1)I(A)$ and
\begin{align}
  U^h-U^g &= (\rho_1-1)\big\{I(A_0)-\mathbb{E}[I(A_0+\varepsilon)]\big\}, \nonumber
\end{align}
positive if $I$ is concave, zero if linear, negative if convex (Jensen, $\rho_1>1$).

\emph{Numerical illustration.} $p_H=0.8$, $\Delta p=0.4$, $R=1.5$ ($\rho_1=1.2$, $p_H/\Delta p=2$),
$B(I)=0.3I+0.1g(I)$, so $\phi(I)=(1-1.2)I+2(0.3I+0.1g)=0.4I+0.2g(I)$; $A_0=10$, $\varepsilon=\pm5$.
\begin{itemize}
  \item Linear ($g=0$): $I=A/0.4$: $I(5)=12.5$, $I(15)=37.5$, mean $25=I(10)$: indifferent.
  \item Convex ($g=I^2$): $0.2I^2+0.4I-A=0$, $I=[-0.4+\sqrt{0.16+0.8A}]/0.4$: $I(5)=4.0990$,
        $I(10)=6.1414$, $I(15)=7.7178$; mean $5.9084<6.1414$: hedge (gain $0.2(6.1414-5.9084)=0.0466$).
  \item Concave ($g=\sqrt I$): with $s=\sqrt I$, $0.4s^2+0.2s-A=0$, $s=[-0.2+\sqrt{0.04+1.6A}]/0.8$:
        $I(5)=10.8528$, $I(10)=22.6219$, $I(15)=34.5606$; mean $22.7067>22.6219$: gamble
        (gain $0.2(22.7067-22.6219)=0.0170$).
\end{itemize}
In all three cases $\phi'(I)>0$ everywhere ($0.4$, $0.4+0.4I$, $0.4+0.1/\sqrt I$).

\answer{(iii)} Hedge if $B$ is convex ($I(A)$ concave), indifferent if $B$ is linear
($I=kA$), gamble if $B$ is concave ($I(A)$ convex).

\step{8}{(iv) Unobservable income} Investors can be promised nothing, so they lend nothing and
$I=A$ (investing all cash is optimal by assumption): $U(A)=R(A)$, concave. Hence
$U^h=R(A_0)>\mathbb{E}[R(A_0+\varepsilon)]=U^g$. Example: $R(I)=2\sqrt I$, $A_0=1$,
$\varepsilon=\pm0.5$: $U^h=2\sqrt1=2$; $U^g=\tfrac12(2\sqrt{0.5}+2\sqrt{1.5})=\tfrac12(1.414214+2.449490)=1.931852<2$.

\answer{(iv)} She hedges: with self-financing and a concave technology, the value of cash is
concave.

\step{9}{(v) Date-0 contracting replicates and improves on gambling} Without hedging, each
$\varepsilon$ leads to an investment $I(\varepsilon)$ and to cash used $A(\varepsilon)\le A_0+\varepsilon$,
with $\mathcal P(I(\varepsilon))\ge I(\varepsilon)-A(\varepsilon)$ in every state ($\mathcal P$: pledgeable income).
Taking expectations,
\begin{align}
  \mathbb{E}[\mathcal P(I)] &\ge \mathbb{E}[I]-\mathbb{E}[A(\varepsilon)] \ge \mathbb{E}[I]-A_0, \nonumber
\end{align}
since $\mathbb{E}[A(\varepsilon)]\le\mathbb{E}[A_0+\varepsilon]=A_0$. Now let her sign at date~0:
investors absorb $\varepsilon$ (full hedging), she draws the investment $\tilde I$ from the same
distribution as $I(\varepsilon)$ but independently of $\varepsilon$, contributes
$\mathbb{E}[A(\varepsilon)]$ and keeps $A_0-\mathbb{E}[A(\varepsilon)]$. Investors' expected profit is
$\mathbb{E}[\mathcal P(\tilde I)]-\mathbb{E}[\tilde I]+\mathbb{E}[A(\varepsilon)]\ge 0$ by the inequality above,
incentives given $\tilde I$ are unchanged, and her expected NPV and leftover cash are the same: she
is at least as well off. She can do strictly better when gambling wastes resources. In case (i)
with $A_0<\bar A$, funding happens with probability $1-F(\bar A-A_0)$ and
$[1-F(\bar A-A_0)]\bar A<A_0$ (cash is unused in every state); a date-0 random-financing contract
(Exercise~5) funds the project with probability $x=A_0/\bar A>1-F(\bar A-A_0)$. In the
illustration with $A_0=30$: $[1-F]\bar A=0.25(40)=10<30$; $x=30/40=0.75$ and
$U=30+0.75(20)=45>35$.

\answer{(v)} With a date-0 contract, any outcome of gambling can be replicated with investment
independent of $\varepsilon$ (full hedging), and often improved on (here $45$ against $35$): the
taste for gambling in (i) and (iii) comes from the absence of contracting at date~0.

\intuition Risk neutrality does not imply indifference to hedging: what matters is the
curvature of the value of internal funds. Financing frictions can make it concave (continuous
effort, self-financed concave technology, convex private benefits), so hedging pays, or convex
(the funding threshold, concave private benefits), so gambling pays. When the firm can arrange
its financing before the shock, it can buy the ``lottery'' it wants directly, without wasting cash,
and insulating investment from the shock becomes optimal.
\pitfalls
\begin{itemize}
  \item Concluding that a risk-neutral entrepreneur never cares about hedging; the curvature of
        $U(A)$, not her utility function, decides.
  \item In (ii), writing the NPV without the effort cost ($R_bR-I$); the conclusion survives
        but $V'$ is then wrong.
  \item Forgetting $\mathbb{E}[\varepsilon]=0$: hedging changes the risk, not the mean of cash on hand.
\end{itemize}
\bookans (i), (iii), (iv) and (v) match the book's answer (pp.~581--582). Two display slips in the
book: in (i) the date-1 gross utility is printed as $[p_HR-I]+[A-\bar A]$ where $[p_HR-I]+A$ is
meant (its next line, $U_b^h=[p_HR-I]+A_0$, uses $A$); in (ii) it prints $U_b=R_bR-I$ and
$dV/dA=R/(2R_b(A)-R)$, omitting the effort cost $\tfrac12R_b^2$; with it,
$dV/dA=(R-R_b)/(2R_b-R)$, as in its own answer to Exercise~3.20 (at the illustration point,
$0.202782$ instead of $1.405564$). Both are still positive with negative second derivative, so
the conclusion (hedge) is unchanged.
\refbib{Tirole Ch.~3, \S3.2.2, pp.~116--117; \S3.4, pp.~127--128; Exercise~3.20, p.~152; Ch.~5, \S5.4 ``Corporate Risk Management'', pp.~213--219.}

\end{document}
